% Author: Joseph Rowell
% Version: 3.0
% This work is licensed under a Creative Commons Attribution 4.0 International License.
\documentclass[12pt, a4paper]{report}
\usepackage[a4paper, total={6.25in, 8.25in}]{geometry} %%%%%% CHECK MARGIN REQ. 8.25 in?
\input{Packages.tex}
\usepackage{multicol}
\hypersetup{pdftitle = Large-Area Satellite-Derived Urban Reconstruction and Manifest-Driven Integration into ChordAtlas, pdfauthor = {Renjie Li}, pdfstartview=FitH, pdfkeywords = urban reconstruction, pdfpagemode = FullScreen, colorlinks, anchorcolor = black, citecolor = blue, urlcolor = blue, filecolor = green, linkcolor = blue, plainpages = false}
%%%%%%%%%%%%%%%%%%%%%%%%%%%%%%%%%%%%%%%%%%%%%%%%%%%%%%%%%%%%%%%%%%%%%%%
\pagestyle{fancy}
\rhead{}
\chead{}
\lhead{University College London}
\lfoot{\date{}}
\cfoot{}
\rfoot{\thepage}
% Top and Bottom Line Rules

\renewcommand{\headrulewidth}{0.4pt} %0.4pt
\renewcommand{\footrulewidth}{0.4pt}
\fancyheadoffset{8pt}
\fancyfootoffset{8pt}
% Line spacing
\renewcommand{\baselinestretch}{1.5} %1.5

\makeglossaries

% Acronyms used in the dissertation
\newacronym{aoi}{AOI}{area of interest}
\newacronym{api}{API}{application programming interface}
\newacronym{bng}{BNG}{British National Grid}
\newacronym{cli}{CLI}{command-line interface}
\newacronym{crs}{CRS}{coordinate reference system}
\newacronym{dsm}{DSM}{digital surface model}
\newacronym{edt}{EDT}{Event Dispatch Thread}
\newacronym{epsg}{EPSG}{European Petroleum Survey Group geodetic registry}
\newacronym{gis}{GIS}{geographic information system}
\newacronym{gui}{GUI}{graphical user interface}
\newacronym{lidar}{LiDAR}{light detection and ranging}
\newacronym{obj}{OBJ}{Wavefront object mesh format}
\newacronym{odbl}{ODbL}{Open Data Commons Open Database Licence}
\newacronym{ogl}{OGL}{Open Government Licence}
\newacronym{osm}{OSM}{OpenStreetMap}
\newacronym{ply}{PLY}{Polygon File Format}
\newacronym{rmse}{RMSE}{root-mean-square error}
\newacronym{wgs84}{WGS84}{World Geodetic System 1984}

% Project-specific terms
\newglossaryentry{manifest}{
  name={manifest},
  description={Machine-readable record that binds a job or publication to its inputs, configuration, state, and outputs}
}

\newglossaryentry{minimesh}{
  name={MiniMesh},
  plural={MiniMeshes},
  description={ChordAtlas tiled mesh representation comprising an index and spatially partitioned model files}
}

\newglossaryentry{coarseready}{
  name={\texttt{COARSE\_READY}},
  description={On-demand publication state indicating that a building passed the configured coarse geometric gates; it does not mean that a complete or ground-truth-accurate building was reconstructed}
}

\newglossaryentry{ready}{
  name={\texttt{READY}},
  description={Publication state written after the Python assets and metadata pass their contract checks; Java MiniMesh conversion and GUI loading occur afterwards}
}

\newglossaryentry{localframe}{
  name={local frame},
  description={Project coordinate convention in metres with X pointing east, Y representing height, and Z pointing south from a recorded WGS84 origin}
}

\newglossaryentry{provenance}{
  name={provenance},
  description={Traceable relationship among source data, code revision, configuration, model, command, time, and generated artefact}
}

\newglossaryentry{failclosed}{
  name={fail-closed},
  description={Behaviour in which an invalid or incomplete prerequisite prevents later processing and prevents publication of a misleading success state}
}


\date{September 2026}

\title{Large-Area Satellite-Derived Urban Reconstruction and Manifest-Driven Integration into ChordAtlas}
\author{\\ \Large{Renjie Li}
\\
\\ Supervisor: Professor Anthony Steed
\\ Faculty of Engineering, Department of Computer Science
\\ University College London\\[0.4cm]
A Project Report Presented in Partial Fulfillment of the Degree \\ \textit{MSc Computer Graphics, Vision and Imaging}
\\[0.6cm]
\parbox{0.88\textwidth}{\centering\small
This report is submitted as part requirement for the MSc Degree in
'Computer Graphics, Vision \& Imaging', at University College London. It is
substantially the result of my own work except where explicitly indicated in
the text.\\[0.4em]
The report will be distributed to the internal and external examiners, but
thereafter may not be copied or distributed except with permission from the
author.}
}
%%%%%%%%%%%%%%%%%%%%%%%%%%%%%%%%%%%%%%%%%%%%%%%%%%%%%%%%%%%%%%%%%%%%%%%
\begin{document}
%Adjust logo positions here
% \AddToShipoutPicture*{\parbox[t][\paperheight][t]{\paperwidth}{%
%           \includegraphics[width=\paperwidth]{\BackgroundPicturea{Logos/ucl_long%_logo.png}{3in}{3in}}
%           }}
% \AddToShipoutPicture*{\centering\BackgroundPictureb{Logos/Bentham2011_065_c623d.jpg}{3in}{3.7in}}
\AddToShipoutPictureBG*{%
  \AtPageUpperLeft{%
    \raisebox{-\height}{%
      
\includegraphics[width=\paperwidth]{Logos/UCL page header.png}%
    }}
}
\AddToShipoutPicture*{%
      \parbox[t][\paperheight][t]{\paperwidth}{%
          
\includegraphics[width=1.2\paperwidth]{Logos/UCL page footer.png}
      }}


\thispagestyle{headings}
\maketitle
\FloatBarrier
\pagenumbering{roman}

\thispagestyle{empty}
\begin{abstract}
Recent AI-based approaches, particularly deep-learning models, can reconstruct three-dimensional urban geometry from satellite imagery. But bounded inference windows, discontinuities between overlapping
predictions, incomplete large-area appearance output, and limited geospatial
interoperability constrain regional use.  This dissertation develops and
evaluates a hybrid workflow that extends Sat3DGen large-image inference and
connects its outputs to open geospatial data and ChordAtlas.

The proposed large-image method processes a rectangular satellite mosaic with
overlapping windows, maps their predictions into a shared density field at
sub-voxel precision, and blends them with raised-cosine weights before a single
surface extraction.  A memory-bounded second pass reconstructs the window
features and blends per-vertex colour on the fixed global mesh.  In a London case study, this reduced mean density discontinuities at the
former window boundaries by 49--60 per cent and mean top-surface
discontinuities by 60--75 per cent relative to a matched hard-crop baseline.
Every extracted vertex received colour support without changing the mesh
geometry.

OpenStreetMap footprints and a one-metre Environment Agency digital surface
model then provide geographic identity, semantic context, and a common
vertical reference.  The resulting assets are cropped by footprint, retain
vertex RGB, and are published through explicit coordinate, manifest, cache,
and validation contracts to the ChordAtlas environment.  A separate
mesh-space study evaluates the behaviour of independent-tile stitching and
DSM correction.  Together, the results establish a large-area reconstruction
and asset-integration framework that combines learned density and appearance
with explicit spatial and publication contracts for editable urban content.


\keywords{urban reconstruction - satellite imagery - large-image inference - density fusion - vertex colour - ChordAtlas - geospatial integration}
% \vspace{-10mm} %To remove added white space after
\end{abstract}
\newpage
\thispagestyle{empty}
\begin{center}
To my family and everyone who supported this work.
\end{center}

\newpage
\thispagestyle{empty}
\vspace*{\fill}
\begin{center}
Copyright \copyright  \thinspace 2026 by Renjie Li \\ All Rights Reserved
\end{center}
\vspace*{\fill}
\newpage
\thispagestyle{empty}
%\epigraph{The Book of Nature is written in the language of mathematics.}{--- \textup{Galileo Galilei}}
\epigraph{The computer was born to solve problems that did not exist before.}{--- \textup{Bill Gates}}


\thispagestyle{empty}
\chapter*{Acknowledgements}

I am grateful to Professor Anthony Steed for his supervision and for shaping
the project from a broad world-generation theme into a concrete investigation
of geographically steered urban reconstruction.


\thispagestyle{empty}
\chapter*{Declaration}
I, Renjie Li, declare that the thesis has been composed by myself and that the work has not been submitted for any other degree or professional qualification. I confirm that the work submitted is my own, except where work which has formed part of jointly-authored publications has been included and referenced. The report may be freely copied and distributed provided the source is explicitly acknowledged. \\
%Add signature png here.
% \begin{figure}[H]
% \includegraphics{Logos/signature.PNG}
% \end{figure}
\vspace{-2cm}
\noindent\begin{tabular}{ll}
 & 14/08/26 \\
\makebox[2.5in]{\hrulefill} & \makebox[2.5in]{\hrulefill}\\
\textit{Signature} & \textit{Date}\\
\end{tabular}



\tableofcontents

\thispagestyle{plain}
\listoffigures
\listoftables
%\lstlistoflistings
\listofalgorithms


%\begin{singlespace}
\chapter*{List of Abbreviations}
\begin{sortedlist} %sort alphabetically
  \sortitem{AOI: Area of Interest}
  \sortitem{API: Application Programming Interface}
  \sortitem{BNG: British National Grid}
  \sortitem{CRS: Coordinate Reference System}
  \sortitem{DSM: Digital Surface Model}
  \sortitem{GIS: Geographic Information System}
  \sortitem{GUI: Graphical User Interface}
  \sortitem{LiDAR: Light Detection and Ranging}
  \sortitem{OBJ: Wavefront Object Mesh Format}
  \sortitem{OSM: OpenStreetMap}
  \sortitem{WGS84: World Geodetic System 1984}
  \end{sortedlist}
% \end{singlespace}
%%%%%%%%%%%%%%%%%%%%%%%%%%%%%%%%%%%%%%%%%%%%%%%%%%%%%%%%%%%%%%%%%%%%%%%%%%%%%%%%
\chapter{Introduction} \label{Chap1}\label{chap:introduction}
\pagenumbering{arabic}



\section{Motivation and problem statement}

Deep generative models conditioned on widely available satellite imagery offer
a promising route to detailed urban geometry.  Their outputs are not,
however, immediately suitable for an editable city-modelling environment.
Inference is performed over bounded image windows; neighbouring predictions
may disagree in position, height, and topology; and a mesh file does not by
itself identify its geographic extent, building ownership, coordinate frame,
or processing history.  These issues become particularly important when the
output must support procedural editing or interactive use rather than a single
rendered view.  Figure~\ref{fig:sat3dgen-tile-example} illustrates the form of
one retained local output.

\begin{figure}[H]
\centering
\includegraphics[width=0.98\textwidth]{sat3dgen_tile_example.pdf}
\caption[Representative Sat3DGen mesh from one London inference tile.]{Representative
Sat3DGen output for a London inference tile centred at $51.510303^{\circ}$~N,
$0.133978^{\circ}$~W.  Panel (a) renders the stored predicted vertex colours;
panel (b) renders the same vertices and triangles with a uniform material.  The
paired views separate predicted appearance from the generated surface in the
model's normalised local frame.}
\label{fig:sat3dgen-tile-example}
\end{figure}

Procedural city systems address a complementary part of the problem.  Building
footprints, semantic parts, and explicit rules provide stable identity and
authoring control, but do not recover the observed appearance and geometric
detail of a real location automatically.  This dissertation studies the
interface between these representations.  It asks how learned urban geometry
can be combined with open geospatial context and delivered as building-level
assets to ChordAtlas, an existing procedural urban-modelling environment
\cite{kellyChordAtlas2024}.

A case study is centred on London and uses Sat3DGen as the
satellite-conditioned generator \cite{qian2026sat3dgen}.  The study extends its large-image workflow with continuous density placement, smooth
overlap weighting, and vertex-colour recovery before connecting the resulting
surface to downstream geospatial processing.  OpenStreetMap (OSM)
provides footprints and contextual semantics, while an Environment Agency
digital surface model (DSM) supplies an external vertical reference
\cite{openstreetmapData,environmentAgencyDSM2022}.  The resulting problem is
both geometric and operational: tiles must be placed and combined coherently,
and every published asset must retain enough spatial and processing metadata
to be interpreted by the downstream system.

\section{Research aim and question}

The aim is to develop and evaluate a hybrid machine-learning/procedural
workflow that reconstructs a contiguous region from overlapping
satellite-conditioned predictions and converts the resulting geometry into
coordinate-explicit, building-associated assets for ChordAtlas.  The study
addresses the following research question:

\begin{quote}
\textbf{To what extent can overlapping satellite-conditioned predictions be
combined into a coloured regional mesh and integrated with geospatial context
in ChordAtlas through a manifest-driven, coordinate-explicit workflow?}
\end{quote}

The question treats integration as more than file conversion.  A useful result
must establish a consistent spatial frame, associate geometry with geographic
features, preserve appearance information, expose processing state, and
support repeatable transfer between Python mesh processing and the Java-based
ChordAtlas environment.

\section{Design rationale and evaluated scope}

The broader project theme concerned geographically steered world generation
for interactive and collaborative environments.  Sat3DGen was selected because
it produces explicit coloured geometry from satellite imagery, making it a
more direct input to an editing system than an image-only generator.  Initial
feasibility work established local model inference and explored Unity/Ubiq
integration.  A London cross-view collection was then prepared to investigate
location-specific data requirements.

Resource analysis showed that reproducing the released Sat3DGen training
regime would dominate the available computational budget.  The empirical study
therefore retains the released model and evaluates two complementary forms of
regional assembly.  The first adapts the released large-image wrapper so that
rectangular mosaics are fused smoothly in density space and colour is recovered
on the global surface.  The second places independently extracted meshes in a
shared geographic frame and evaluates stitching, height reconciliation, and
building association.  Both paths subsequently enter the same publication
architecture.

OSM and DSM data serve different roles in this design.  Footprints provide
building identity and a basis for per-building extraction; the DSM provides a
metric surface-height reference for vertical processing.  ChordAtlas provides
the procedural destination, including its tiled MiniMesh representation and
building-generation workflow.  BigSUR and FrankenGAN inform the wider goal of
structured fa\c{c}ade modelling, while CityEngine provides a procedural
workflow comparator and Ubiq motivates future collaborative deployment
\cite{kelly2017bigsur,kelly2018frankengan,esriCityEngineIntro2026,
friston2021ubiq}.  These systems establish the architectural context for the
study; the evaluated contribution comprises the London data study, large-image
density and colour method, mesh-processing experiments, and ChordAtlas
integration path.

\section{Objectives}

Six objectives make the research question operational:

\begin{enumerate}
  \item construct and characterise a London multimodal case study comprising
  satellite imagery, panoramas, OSM context, and DSM data;
  \item adapt and evaluate large-image density fusion and per-vertex colour
  recovery while retaining the released Sat3DGen model;
  \item define the coordinate, tiling, and metadata contracts that connect
  geographic inputs, Sat3DGen density and mesh spaces, and ChordAtlas;
  \item implement mesh-space placement, cropping, overlap processing,
  DSM-informed height correction, and footprint-based building extraction;
  \item implement static and on-demand publication of building assets to
  ChordAtlas while preserving stable identity and vertex colour; and
  \item evaluate data integrity, geometric behaviour, topology, coordinate
  consistency, and cross-language publication, and relate the findings to
  Sat3DGen and procedural city-generation methods.
\end{enumerate}

\section{Contributions}

The first contribution is a London data-preparation and integrity study that
connects cross-view imagery with OSM and DSM context.  It quantifies geographic
support, channel availability, duplicate imagery, split relationships, and
candidate-label geometry.  These measurements show which properties must be
checked before a geographically assembled collection is used for learning or
evaluation.

The second contribution is a rectangular large-image inference method.  It
places overlapping density predictions at sub-voxel coordinates, blends them
with smooth spatial weights, extracts one global surface, and recovers
per-vertex RGB in a memory-bounded second pass.  A matched experiment measures
the reduction in window-boundary discontinuity relative to hard-crop fusion.

The third contribution is a mesh-space integration method for placing and
combining independently generated tiles.  The method introduces explicit
coordinate conversion, footprint association, DSM-informed vertical
processing, and per-building extraction.  Controlled geometry cases and
topological measurements evaluate the assumptions made by bottom removal,
stitching, and height correction.

The fourth contribution is a ChordAtlas bridge that publishes selected geometry
as modular building assets.  Stable request identities, manifests, validation
gates, staged publication, vertex-colour transfer, and Python--Java conversion
make the data boundary inspectable and repeatable.  The evaluation identifies
the conditions under which this bridge operates reliably and the geometry and
transaction properties that most strongly affect regional extension.

Sat3DGen, ChordAtlas, BigSUR, FrankenGAN, CityEngine, and Ubiq are pre-existing
systems.  The project-specific work reported here comprises the London data
workflow, large-image density and colour extension, geospatial mesh-processing
method, ChordAtlas integration, and their evaluation.  External software,
datasets, and published methods are identified and cited where they are used.

\section{Dissertation structure}

Chapter~\ref{chap:background} reviews procedural urban modelling,
satellite-conditioned reconstruction, semantic city representations,
geospatial contracts, and related systems.  Chapter~\ref{chap:methodology}
defines the data, large-image inference, coordinate, mesh-processing,
publication, and evaluation methods.  Chapter~\ref{chap:results} reports the
data-integrity, density-fusion, colour, geometry, coordinate, and ChordAtlas
integration results.  Chapter~\ref{chap:discussion}
interprets the findings, compares the approach with Sat3DGen and CityEngine,
and evaluates validity and future directions.  Chapter~\ref{chap:conclusion}
summarises the research outcome and contributions.  The appendix records
implementation provenance and detailed component-test results.

%%%%%%%%%%%%%%%%%%%%%%%%%%%%%%%%%%%%%%%%%%%%%%%%%%%%%%%%%%%%%%%%%%%%%%%%%%%%%%%%
\chapter{Background and Related Work} \label{Chap2}\label{chap:background}

\section{Procedural and data-driven urban modelling}

Procedural modelling represents architecture through rules and parameters
rather than solely through sampled surfaces.  Early city systems used extended
L-systems to generate street networks and building populations
\cite{parish2001procedural}; split grammars then showed how attributes and
rules could control architectural variation at scale \cite{wonka2003instant}.
CGA shape grammar made the relationship between initial shapes, rule
operations, and polygonal buildings explicit \cite{muller2006procedural}.
ProceX made procedural extrusions
interactive: footprint-driven operations create architectural massing that can
be edited rather than merely viewed \cite{kelly2011procex}.  This structured
representation is attractive for planning, simulation, visualisation, and
content authoring because it preserves semantic parts.

Open building-footprint collections provide initial shapes for applying this
principle beyond manually authored scenes.  Microsoft's Global ML Building
Footprints dataset supplies predicted building outlines at large geographic
scale \cite{microsoftGlobalFootprints2026}.  These outlines provide location
and plan extent, while height, appearance, and architectural semantics must
still be reconstructed or generated.

Real data rarely arrives as clean procedural parameters.  BigSUR addresses a
closely related integration problem by combining noisy large-scale meshes,
imagery, and \gls{gis} footprints to recover structured urban models
\cite{kelly2017bigsur}.  FrankenGAN then enriches coarse mass models with
semantic fa\c{c}ade detail and appearance through style-synchronised generative
networks \cite{kelly2018frankengan}.  ChordAtlas exposes these research
components in an interactive Java environment, with generators for GIS data,
panoramas, meshes, and blocks \cite{kellyChordAtlas2024}.  FrankenGAN assumes a
structured building-mass representation: it is neither a topology-repair
method for arbitrary meshes nor a source of ground-truth fa\c{c}ade texture.
ChordAtlas appearance exports are therefore treated here as procedurally
generated assets rather than texture-preserving conversions of Sat3DGen
triangles.

Recent geospatial approaches divide geometry and appearance processing in
different ways.  GeoTexBuild begins with map footprints and separates
height-map generation, geometry reconstruction, and appearance stylisation
\cite{wang2025geotexbuild}.  Buyukdemircioglu and Oude Elberink instead texture
existing CityJSON geometry from oblique and nadir aerial imagery
\cite{buyukdemircioglu2024texture}.  Their contrasting starting points clarify
that geometric reconstruction, structured city representation, and appearance
transfer are related but distinct stages.

ArcGIS CityEngine is a relevant industrial authoring system in this procedural
lineage.  Its central workflow combines streets, blocks, lots, initial shapes,
GIS attributes, assets, and authored CGA rules to generate polygonal city
content \cite{esriCityEngineIntro2026}.  File-geodatabase attributes can be
retained and used to drive rules, while the documented runtime and SDK support
scripted generation and export \cite{esriCityEngineGIS2026,
esriCityEngineSDK2026}.  These capabilities favour controllable regeneration,
scenario authoring, and attribute-aware content production.  They do not,
however, imply automatic recovery of a real building from satellite pixels.
Likewise, importing a dense OBJ does not by itself create stable building
identity or roof, wall, and window semantics.  CityEngine therefore serves as
a representation and workflow comparator, including an executed local
example, but not as a satellite-reconstruction accuracy baseline.

This lineage establishes two requirements for the present integration.  First,
generated geometry must be spatially compatible with footprints and panoramas,
not merely visually plausible in isolation.  Secondly, it must be converted
into ChordAtlas's tiled \glspl{minimesh} and generator graph without corrupting
an existing project when conversion fails.  The baseline ChordAtlas workflow
can clip an already loaded mesh and construct a block.  The proposed path adds
external, on-demand generation and publication, so state and failure handling
become part of the research problem.

\section{Satellite-conditioned street-level geometry}

Satellite-to-ground synthesis is underconstrained: roofs and plan-view context
are visible from above, while fa\c{c}ades and street-level occlusions are not.
Cross-view datasets reduce that ambiguity by pairing overhead imagery with
ground panoramas.  VIGOR formalised a setting in which a query is not
necessarily centred in its reference satellite image and uses overlapping
reference tiles \cite{zhu2021vigor}.  Its correspondence geometry is therefore
not an incidental file format.  Candidate offsets must describe where the
ground query lies within each aerial tile, and sufficient overlap is required
if several candidates are meant to contain the same location.

Sat2Density learns a neural density representation from satellite--ground
pairs \cite{qian2023sat2density}.  Its successor extends the cross-view
rendering lineage \cite{qian2026sat2densitypp}.  Sat2Scene uses diffusion for urban scene
generation \cite{li2024sat2scene}, while Sat2City uses a cascaded latent
diffusion formulation for city generation \cite{hua2025sat2city}.  These
methods are related baselines, but their tasks, training data, and output
representations differ.  The comparison therefore considers their
architectures and evaluation protocols rather than treating their reported
metrics as interchangeable.

At the broader systems level, Meta's WorldGen and World Labs illustrate a move
towards prompt- and observation-conditioned 3D environments
\cite{metaWorldGen2025,worldLabs2026}.  Their relevance here is the objective of
producing navigable and editable world content from location-specific
observations; they provide systems context rather than directly comparable
Sat3DGen baselines.

Sat3DGen places particular emphasis on explicit geometry.  Its published
pipeline uses a frozen DINOv3 image encoder, spatial padding tokens, a decoder
that produces three feature planes, and a small network that predicts density
and colour \cite{simeoni2025dinov3,qian2026sat3dgen}.  The tri-plane design is
related to efficient geometry-aware neural rendering \cite{chan2022eg3d}.
Training adds a gravity-based density-variation loss, a relative-depth prior
derived from Depth Anything V2 \cite{yang2024depthanythingv2}, and randomly
projected perspective views from ground panoramas.  The density representation
can be rendered from satellite or street-level viewpoints and converted into
an explicit surface through Marching Cubes
\cite{qian2026sat3dgen,lorensen1987marching}.  Figure~\ref{fig:sat3dgen-model}
separates this published model from the large-image and geospatial extensions
studied here.

\begin{figure}[H]
\centering
\resizebox{0.98\textwidth}{!}{%
\begin{tikzpicture}[
  node distance=7mm and 6mm,
  model/.style={draw=evidenceblue,rounded corners,align=center,
    minimum height=11mm,text width=25mm,fill=evidenceblue!8},
  train/.style={draw=evidenceorange,rounded corners,align=center,
    minimum height=9mm,text width=25mm,fill=evidenceorange!8},
  downstream/.style={draw=evidencegreen,rounded corners,align=center,
    minimum height=11mm,text width=32mm,fill=evidencegreen!8},
  arrow/.style={-{Latex[length=2mm]},thick}
]
\node[model] (image) {Satellite RGB\\$256\!\times\!256$};
\node[model,right=of image] (dino) {Frozen DINOv3\\token encoder};
\node[model,right=of dino] (tokens) {$16\!\times\!16$ tokens\\+ spatial padding};
\node[model,right=of tokens] (planes) {Decoder\\three feature planes};
\node[model,right=of planes] (field) {MLP density\\and colour field};
\node[model,right=of field] (mesh) {Rendering or\\Marching Cubes};
\draw[arrow] (image) -- (dino);
\draw[arrow] (dino) -- (tokens);
\draw[arrow] (tokens) -- (planes);
\draw[arrow] (planes) -- (field);
\draw[arrow] (field) -- (mesh);
\node[train,below=11mm of dino] (gravity) {Gravity density\\variation loss};
\node[train,right=of gravity] (depth) {Relative-depth\\prior};
\node[train,right=of depth] (persp) {Perspective\\views from\\panoramas};
\draw[arrow,dashed,evidenceorange] (gravity.north) -- (tokens.south);
\draw[arrow,dashed,evidenceorange] (depth.north) -- (planes.south);
\draw[arrow,dashed,evidenceorange] (persp.north) -- (field.south);
\node[downstream,below=11mm of field] (large) {This dissertation:\\fractional density fusion\\and vertex colour};
\node[downstream,right=16mm of large] (bridge) {Geospatial mesh and\\ChordAtlas bridge};
\draw[arrow,evidencegreen] (field.south) --
  node[midway,right=1mm,font=\scriptsize,align=left]{overlapping\\windows}
  (large.north);
\draw[arrow,evidencegreen] (large.east) --
  node[midway,above=1mm,font=\scriptsize,align=center]{global\\coloured mesh}
  (bridge.west);
\draw[arrow,dashed,evidencegreen] (mesh.south) to[out=-90,in=90]
  node[pos=.55,right=1mm,font=\scriptsize,align=left]{independent\\mesh path}
  (bridge.north);
\end{tikzpicture}%
}
\caption{Published Sat3DGen components and the two project extensions.  Blue
and orange identify the model and training components of Qian et al.; green
identifies the large-image density/colour method and the subsequent geospatial
integration developed in this dissertation.  The dashed link is the separate
independent-mesh experiment \cite{qian2026sat3dgen}.}
\label{fig:sat3dgen-model}
\end{figure}

The paper reports training on VIGOR-derived cross-view pairs, evaluation
outside the training cities, and a large-scale strategy based on overlapping
patches.  Its analysis identifies pose accuracy, orthographic projection,
locally flat ground, and atypical architecture as important reconstruction
factors.  The present study extends the systems context by examining
local-to-world transforms, seam handling, DSM alignment, building association,
and ChordAtlas publication \cite{qian2026sat3dgen}.

Three processing domains must be distinguished.  The released Sat3DGen
large-image wrapper predicts overlapping density volumes, applies hard boundary
cropping and equal-weight averaging, and runs Marching Cubes once on the fused
field; unlike its single-image mesh path, the wrapper does not query or export
vertex colour \cite{qian2026sat3dgen,qianSat3DGenCode2026}.  The large-image
extension developed here retains one global density-domain surface extraction
but maps window locations at sub-voxel precision, feathers overlapping
contributions, and performs a memory-bounded second pass for vertex RGB.  The
subsequent geospatial path operates in mesh space, adding OSM/DSM alignment,
footprint association, per-building extraction, and ChordAtlas publication.

\section{Semantic city models, mesh validity, and provenance}

A renderable triangle soup is not yet a semantic city model.  CityGML 3.0
separates georeferenced geometry from urban-object semantics, relationships,
and multiple spatial representations \cite{ogcCityGML3}.  Its Levels of Detail
should not be reduced to a generic low/medium/high polygon-count switch:
geometric complexity and semantic granularity are distinct
\cite{biljecki2016lod}.  The bridge's footprint-linked identifiers are therefore a
useful coarse association; the current OBJ outputs do not encode the validated
roof and wall semantics required for a CityGML representation.

Downstream validity is similarly task-dependent.  Local mesh repair may close
holes or remove defects, but it can also change the observed surface
\cite{attene2010repair}; standards-oriented validation distinguishes valid
primitives from merely readable files \cite{ledoux2018val3dity}.  This motivates
reporting boundary and non-manifold edges separately from geometric accuracy.

Finally, a manifest is only a container unless it records causal lineage.
PROV-O distinguishes entities, activities, and agents \cite{w3cProvO2013}; in
this setting the equivalent minimum is an exact input allow-list and hashes,
geographic extent and CRS, code/model/configuration identity and execution state.  

\section{Geospatial contracts}

\subsection{Coordinate order and reference systems}

GeoJSON positions are written in longitude--latitude order
\cite{butler2016geojson}, whereas an Overpass QL bounding box is ordered south,
west, north, east \cite{overpassQL}.  Confusing the two conventions can send a
valid query to the wrong region without producing a type error.  This is a
representative integration defect: each subsystem may accept four floating
point values while disagreeing about their meaning.

The data-builder grid uses Web Mercator world pixels at zoom 20.  EPSG:3857 is
appropriate for web mapping, but its registry explicitly notes scale
distortion \cite{epsg3857}.  Metric analysis of the London DSM instead uses
EPSG:27700, the British National Grid \cite{epsg27700}.  The bridge establishes
a local frame whose origin is stored in WGS84 and whose axes are X east, Y
height, and Z south.  The frame is convenient for ChordAtlas and Sat3DGen mesh
operations, but a constant-metres-per-degree approximation is not equivalent
to a projected CRS.  Chapter~\ref{chap:results} quantifies this difference.

\subsection{Surface elevation and semantics}

The Environment Agency's \SI{1}{\metre} LiDAR Composite DSM records surface
height rather than bare-earth terrain.  It can therefore include buildings,
vegetation, and vehicles; its 2022 composite also joins surveys acquired at
different times \cite{environmentAgencyDSM2022}.  GeoTIFF provides the
georeferenced raster interchange contract \cite{ogc2019geotiff}.  A DSM can
help place generated geometry vertically, but it is not an exact and
contemporaneous building-height ground truth.

OSM supplies footprints and contextual categories.  It is a collaborative
database under the ODbL \cite{openstreetmapData}.  Overpass queries commonly
return ways, relations, and the supporting nodes needed to reconstruct their
geometry.  Those support nodes must not automatically become category features,
and multipolygon outer/inner rings must retain their roles.  The data audit in
Chapter~\ref{chap:results} shows why semantic masks cannot be accepted solely
because image files were produced.

\section{Synthesis and research gap}

The literature reviewed above addresses complementary parts of the target
workflow.  Cross-view models connect satellite imagery to density fields,
rendered views, or explicit meshes
\cite{qian2023sat2density,qian2026sat3dgen}; procedural systems support
structured and editable urban content
\cite{kelly2017bigsur,kellyChordAtlas2024,esriCityEngineIntro2026}; and OSM,
DSM, and city-model standards provide geographic context, elevation, and
semantic conventions
\cite{openstreetmapData,environmentAgencyDSM2022,ogcCityGML3}.  These
components use different spatial frames and representations, so their outputs
do not compose automatically.

Among the reviewed approaches, none follows overlapping
satellite-conditioned predictions through regional fusion, geospatial
registration, footprint-based building extraction, and publication into an
interactive procedural environment.  Neural reconstruction work primarily
evaluates geometry or rendering, whereas procedural workflows generally begin
with geometry and geographic structure already available.  The research gap
is therefore the transition between these representations and the evaluation
of that transition as an integrated workflow.

Accordingly, this study evaluates whether satellite tiles, projected DSM data,
local mesh coordinates, and ChordAtlas placement remain spatially consistent.
It measures the effects of fusion, cropping, stitching, vertical correction,
and building extraction.  These evaluation
dimensions are operationalised in Chapter~\ref{chap:methodology} and provide
the basis for the research question introduced in
Chapter~\ref{chap:introduction}.

%%%%%%%%%%%%%%%%%%%%%%%%%%%%%%%%%%%%%%%%%%%%%%%%%%%%%%%%%%%%%%%%%%%%%%%%%%%%%%%%
\chapter{Methodology} \label{Chap3}\label{chap:methodology}

\section{Research scope and methodological rationale}

The empirical design evaluates six linked components: London data preparation,
large-image density and colour inference, independent-mesh processing, static
and on-demand ChordAtlas integration, a matched representation comparison, and
component verification.  Together,
they test whether overlapping satellite-conditioned predictions can form a
continuous coloured surface and be converted into coordinate-explicit,
building-level assets for an interactive procedural environment.

A contiguous Sat3DGen experiment measures overlap fusion and global colour
recovery.  A separate experiment uses retained OBJ outputs to isolate
geospatial placement, mesh stitching, OSM/DSM processing, and building
extraction.  Both routes are evaluated at the ChordAtlas publication boundary.
Data-quality tests characterise the London multimodal collection.  Retained
ChordAtlas appearance exports provide representation-level evidence.

\section{Systems evaluation design}

The evaluation combines data-integrity analysis, coordinate tests, staged
geometry measurements, controlled regression cases, topological inspection,
and cross-language integration tests.  Source tracing establishes the
transformations applied by the active pipeline; generated artefacts and
manifests provide stage-level measurements; and executable checks test the
coordinate, geometry, cache, and publication contracts.  The ChordAtlas bridge
is compared with the existing workflow at the level of representation,
selection, conversion, and publication behaviour.

The research was organised around five evaluation questions:

\begin{description}
  \item[EQ1] Are the source data and pairing conventions internally consistent
  enough to support the intended cross-view and geospatial uses?
  \item[EQ2] How do density-coordinate placement, overlap weighting, and
  second-pass colour recovery affect boundary continuity, colour completeness,
  and resource use in large-image inference?
  \item[EQ3] How did the executed mesh-processing stages affect mesh structure, continuity, and height, and was their observed behaviour consistent with the intended design?
  \item[EQ4] What do the static and on-demand paths demonstrably achieve in
  ChordAtlas, and where do their atomicity and cache boundaries lie?
  \item[EQ5] How do the retained CityEngine and ChordAtlas artefacts differ in
  representation from the two matched Sat3DGen routes, and which comparisons
  remain valid without geometric or appearance ground truth?
\end{description}

\section{Evaluated configuration}

The evaluated configuration comprises the London data-preparation workflow,
the released Sat3DGen checkpoint and large-image extension, the active
independent-mesh package, and the ChordAtlas integration layer.  Runtime imports
and entry-point tracing identify the components used by each experiment.
Functional analysis covers GUI entry points, worker and process boundaries,
stable selection identities, manifest fields, cache validation, publication
order, rollback, and validation depth.  Reproduction details are recorded in
Appendix~\ref{app:evidence-record}.

\section{System overview}

Figure~\ref{fig:system-overview} presents the conceptual data flow derived from
the active implementation.  Chapter~\ref{chap:results} assigns execution and
validation status to the individual paths.

\begin{figure}[H]
\centering
\resizebox{0.98\textwidth}{!}{%
\begin{tikzpicture}[
  x=1mm,
  y=1mm,
  box/.style={draw=uclblue,rounded corners,align=center,minimum width=27mm,
    minimum height=11mm,text width=24mm,fill=lightgrey},
  data/.style={draw=evidencegreen,rounded corners,align=center,minimum width=25mm,
    minimum height=10mm,text width=22mm,fill=evidencegreen!8},
  context/.style={draw=evidencegreen,rounded corners,align=center,minimum width=30mm,
    minimum height=10mm,text width=27mm,fill=evidencegreen!8},
  arrow/.style={-{Latex[length=2.3mm]},thick,uclblue}
]
\node[data] (builder) at (0,22) {London imagery\\and metadata};
\node[box] (sat) at (31,22) {Overlapping Sat3DGen\\window inference};
\node[box] (density) at (63,22) {Fractional density\\fusion and one MC};
\node[box] (colour) at (95,22) {Second-pass\\vertex RGB fusion};

\node[data] (cached) at (0,0) {Cached Sat3DGen\\OBJ tiles};
\node[box] (place) at (31,0) {Coordinate placement\\and crop};
\node[box] (stitch) at (63,0) {Stitching and\\mesh cleanup};
\node[box] (publish) at (95,0) {Building extraction\\and publication};
\node[box] (chord) at (127,0) {ChordAtlas\\renderers and MiniMesh};
\node[context] (geo) at (95,-22) {OSM footprints and\\EPSG:27700 DSM};

\draw[arrow] (builder) -- (sat);
\draw[arrow] (sat) -- (density);
\draw[arrow] (density) -- (colour);
\draw[arrow] (colour) -- (publish);
\draw[arrow] (cached) -- (place);
\draw[arrow] (place) -- (stitch);
\draw[arrow] (stitch) -- (publish);
\draw[arrow] (geo) -- (publish);
\draw[arrow] (publish) -- (chord);
\end{tikzpicture}%
}
\caption[Conceptual architecture of the two evaluated processing routes.]{Conceptual project architecture.  The upper route performs contiguous
density-field fusion and colour recovery; the lower route evaluates
independent-mesh assembly.  Both converge on building extraction and
publication before ChordAtlas.  OSM footprints and the DSM enter from below as
independent geospatial inputs rather than sequential processing stages.}
\label{fig:system-overview}
\end{figure}

\section{London data preparation}

The London data builder is evaluated as a staged data contract supporting
cross-view pairing and downstream geospatial processing.  The protocol measures
coverage, duplication, partition leakage, label geometry, and auxiliary-channel
consistency.
Figure~\ref{fig:data-builder-flow} links each processing stage to its available
artefacts and audit denominator.

\begin{figure}[H]
\centering
\resizebox{0.98\textwidth}{!}{%
\begin{tikzpicture}[
  node distance=8mm and 7mm,
  source/.style={draw=evidenceblue,rounded corners,align=center,
    minimum height=11mm,text width=29mm,fill=evidenceblue!8},
  stage/.style={draw=uclblue,rounded corners,align=center,
    minimum height=11mm,text width=29mm,fill=lightgrey},
  warn/.style={draw=evidenceorange,rounded corners,align=center,
    minimum height=11mm,text width=31mm,fill=evidenceorange!8},
  arrow/.style={-{Latex[length=2mm]},thick,uclblue}
]
\node[source] (points) {OSM-highway\\query points};
\node[stage,right=of points] (match) {Panorama\\matching};
\node[stage,right=of match] (legacy) {Latitude-based\\data splits};
\node[stage,right=of legacy] (grid) {Web Mercator\\tile grid};
\node[stage,right=of grid] (overlap) {Overlapping\\grid variants};
\draw[arrow] (points) -- (match);
\draw[arrow] (match) -- (legacy);
\draw[arrow] (match) -- (grid);
\draw[arrow] (grid) -- (overlap);

\node[source,below=15mm of match] (pano) {Panorama branches\\legacy + equirectangular};
\node[warn,right=of pano] (labels) {VIGOR-like labels\\4 candidates/query};
\node[warn,right=of labels] (channels) {Auxiliary channels\\DSM / OSM\\sky / relative depth};
\node[source,right=of channels] (layouts) {London and Sat3DGen\\directory layouts};
\draw[arrow] (match) -- (pano);
\draw[arrow] (grid) -- (labels);
\draw[arrow] (grid) -- (channels);
\draw[arrow] (overlap) |- (layouts);
\draw[arrow] (pano) -- (labels);
\draw[arrow] (labels) -- (channels);
\draw[arrow] (channels) -- (layouts);
\end{tikzpicture}%
}
\caption{Data-builder stages and the contracts tested for completeness and
geometric consistency.}
\label{fig:data-builder-flow}
\end{figure}

\subsection{Point matching, tiling, and imagery}

The data-builder begins from WGS84 point features derived from an OSM highway
query.  A manually defined longitude--latitude region and a zoom-20 Web
Mercator grid assign satellite centres.  A panorama is matched when its point
lies inside a satellite tile's half-width and half-height.  Chapter~\ref{chap:results}
reports the resulting population and coverage.

For longitude $\lambda$ in degrees, latitude $\phi$ in radians, zoom $z$, and
world-pixel scale $S=256\,2^z$, the scripts use the standard Web Mercator form

\begin{align}
p_x &= S\frac{\lambda+180}{360},\\
p_y &= S\left[\frac{1}{2}-\frac{1}{4\pi}
  \ln\left(\frac{1+\sin\phi}{1-\sin\phi}\right)\right].
\end{align}

Latitude is clamped to the Web Mercator limit.  The nominal ground resolution
at latitude $\phi$ is

\begin{equation}
r(\phi,z)=\frac{2\pi R\cos\phi}{256\,2^z},
\end{equation}

where $R$ is the Web Mercator sphere radius.  At $\phi=51.5074^\circ$ and
$z=20$, the implementation gives approximately
\SI{0.093}{\metre\per\pixel}, so a 640-pixel tile spans about
\SI{59.5}{\metre}.  Values used in the audit were recalculated from the
coordinate metadata.

The acquisition procedure requests 640-by-320 Street View images with a
120-degree field of view and 640-by-640 satellite images, nominally at zoom 20.
The legacy branch does not retain the returned panorama identifier or snapped
coordinates.  The equirectangular branch requests six cube faces and projects them into 1,536-by-768
equirectangular panoramas.  The projection and image dimensions were treated
as data transformations, while observation independence was tested separately
through content hashes.

\subsection{Splits and candidate labels}

The legacy training, validation, and test split sorts by latitude and assigns
70, 15, and 15 per cent respectively.  Spatial sorting can reduce arbitrary
mixing, but there is no boundary buffer, and duplicate returned panoramas can
cross a boundary.

The VIGOR-style branch fills a 16-by-17 Cartesian grid of observed
latitude and longitude centres and chooses the four nearest satellite centres
for every panorama.  For candidate $j$, an offset was classified as containing
the query only when both absolute pixel components were at most 320.  The
audit recorded the count outside this square by candidate rank and the maximum
absolute component.  This tests geometric compatibility rather than mere
four-column syntactic compatibility with VIGOR \cite{zhu2021vigor}.

Two denser satellite grids were also reconstructed.  A 320-pixel stride gives
approximately 50 per cent overlap, while a 576-pixel stride gives 10 per cent
overlap.  Coverage calculations account for coordinate rounding and incomplete
placement at non-divisible outer boundaries.

\subsection{Auxiliary channels}

The pipeline includes binary sky masks, model-derived relative depth, a
DSM-derived normalised raster, and OSM semantic layers.  The DSM branch
transforms WGS84 coordinates into the raster CRS, reprojects a patch with
bilinear sampling, and scales its 5th--95th percentile range independently.
Consequently, it was evaluated as a normalised surface-height channel, not
metric depth.

For OSM, raw element types and tags were counted separately from raster output.
Supporting nodes returned by recursive Overpass queries were identified by
their lack of category tags.  Multipolygon reconstruction and inner-ring
assignment were traced through geometry construction and drawing.  This made
it possible to distinguish an existing image from a semantically valid mask.

\section{Coordinate and vertical-frame contract}

The bridge defines a \gls{localframe} at $(\lambda_0,\phi_0)$ using
$M=111{,}320$ metres per latitude degree:

\begin{align}
x &= (\lambda-\lambda_0)M\cos\phi_0,\\
z &= -(\phi-\phi_0)M.
\end{align}

Thus X points east, Z points south, and Y is vertical.  The inverse uses the
same constants, so algebraic round trips should be exact to floating-point
precision.  The frame is illustrated in Figure~\ref{fig:coordinate-frame}.

\begin{figure}[H]
\centering
\begin{tikzpicture}[scale=1.0,>=Latex]
  \coordinate (o) at (0,0);
  \draw[->,very thick,evidenceblue] (o) -- (4.2,0) node[right]{X: east};
  \draw[->,very thick,evidencegreen] (o) -- (0,3.0) node[above]{Y: height};
  \draw[->,very thick,evidenceorange] (o) -- (-2.2,-1.5) node[below left]{Z: south};
  \fill (o) circle (2.2pt) node[above right]{stored WGS84 origin};
  \draw[dashed,gray] (-3.0,0) -- (4.8,0);
  \draw[dashed,gray] (0,-2.0) -- (0,3.5);
  \node[align=left,anchor=west] at (4.8,2.0)
    {$x=(\lambda-\lambda_0)M\cos\phi_0$\\
     $z=-(\phi-\phi_0)M$};
\end{tikzpicture}
\caption{Local coordinate contract used by the bridge.  The reversible local
approximation is convenient for block-scale processing but is not identical to
EPSG:27700.}
\label{fig:coordinate-frame}
\end{figure}

Coordinate evaluation had two parts.  First, the AOI centre and four corners
were transformed to local coordinates and back, measuring residual in metres.
Secondly, the same points were transformed from EPSG:4326 to EPSG:27700 using
\code{pyproj}; projected easting/northing differences from the origin were
compared with local X/south values.  This separates inverse consistency from
external metric agreement.

The vertical contract has three distinct quantities: generated-mesh local Y,
DSM surface elevation, and a workspace vertical offset.  Their values were
read from manifests and paired OBJ vertices.

\section{Large-image and mesh-processing protocols}

\subsection{Rectangular density-field fusion}

The large-image extension retains the released Sat3DGen checkpoint and its
standard 256-by-256 network transform.  It accepts a rectangular source mosaic
without first truncating it to a central square.  In the evaluated setting, a
640-pixel source window moves by 160 pixels, giving 75 per cent nominal overlap;
each source window is resized by the unchanged model transform.  The
\code{mesh resolution} parameter is the volumetric sampling size, not the image
input size.  A value of $N=192$ and the checkpoint's 0.8 position scale produce
a cropped density patch of $154\times154\times173$ samples with a 19-voxel
border.

Let $s_i$ be the source-image start of window $i$ on either horizontal axis and
$L=154$ the cropped density width.  Its continuous density origin is

\begin{equation}
o_i=s_i\frac{L}{640}.
\end{equation}

The origin is retained as a floating-point value rather than rounded to an
integer cell.  Each density sample is bilinearly distributed to the four
neighbouring cells of the global horizontal grid.  A separable raised-cosine
window tapers the interior 19-voxel boundary of each patch; an edge coincident
with the outer mosaic boundary retains unit weight.  For a global density
position $p$, the accumulated field is

\begin{equation}
\bar{D}(p)=\frac{\sum_i w_i(p)D_i(p)}{\sum_i w_i(p)},
\label{eq:density-feather-fusion}
\end{equation}

where $D_i$ is the resampled prediction and $w_i$ combines its bilinear
placement and spatial taper.  A zero denominator is treated as an invalid
coverage state.  Marching Cubes is applied once to $\bar{D}$ at the evaluated
isovalue of 4.5, so overlap is resolved before surface extraction.

Algorithm~\ref{alg:fractional-density-fusion} gives the corresponding
semi-formal procedure.  The two accumulators separate predicted density from
its support weight, which permits a final normalised field to be extracted by
a single Marching Cubes operation.

\begin{algorithm}[H]
\caption{Fractional feather fusion of overlapping density predictions}
\label{alg:fractional-density-fusion}
\begin{algorithmic}[1]
\Require Rectangular mosaic $I$; source-window starts $S$; predictor $M$;
cropped density width $L$; halo $h$; isovalue $\tau$
\Ensure Global vertices $V$ and faces $F$
\State Initialise density sum $A\gets 0$ and horizontal weight sum $W\gets 0$
\ForAll{$(r_i,c_i)\in S$}
  \State $D_i\gets\Call{PredictDensity}{M,I,r_i,c_i}$
  \State $D_i\gets\Call{CropHalo}{D_i,h}$
  \State $(o_{r},o_{c})\gets(r_iL/640,c_iL/640)$
  \State $K_i\gets\Call{RaisedCosineWeights}{L,h,\text{mosaic-edge flags}}$
  \State $(f_r,f_c)\gets(o_r-\lfloor o_r\rfloor,o_c-\lfloor o_c\rfloor)$
  \ForAll{$(a,b)\in\{0,1\}^{2}$}
    \State $\alpha_{ab}\gets(af_r+(1-a)(1-f_r))(bf_c+(1-b)(1-f_c))$
    \State \Call{Accumulate}{$A,\alpha_{ab}K_iD_i,\lfloor o_r\rfloor+a,
      \lfloor o_c\rfloor+b$}
    \State \Call{Accumulate}{$W,\alpha_{ab}K_i,\lfloor o_r\rfloor+a,
      \lfloor o_c\rfloor+b$}
  \EndFor
\EndFor
\State \textbf{require} $W(p)>0$ for every covered horizontal position $p$
\State $\bar D\gets A/W$ \Comment{broadcast $W$ over the vertical axis}
\State $(V,F)\gets\Call{MarchingCubes}{\bar D,\tau}$
\State \Return $(V,F)$
\end{algorithmic}
\end{algorithm}

\subsection{Memory-bounded vertex-colour recovery}

The released large-image wrapper discards each tri-plane after density
prediction and therefore exports geometry without the vertex-colour query used
by single-image inference \cite{qianSat3DGenCode2026}.  The added colour pass
keeps the extracted global vertices and faces fixed.  It reprocesses one source
window at a time, reconstructs that window's tri-plane, identifies the global
vertices supported by the same fractional feather weights, and queries the
pretrained colour network in batches.

The global Marching-Cubes axes are $[\mathrm{row},\mathrm{column},z]$.  For a
window with continuous origin $(o_r,o_c)$, a supported vertex is mapped back to
the full $192^3$ model grid as

\begin{equation}
q=\left[\mathrm{column}-o_c+19,
         \mathrm{row}-o_r+19,
         z+19\right]\frac{2}{191}-1.
\label{eq:large-image-colour-coordinate}
\end{equation}

The horizontal query is clamped only for the fractional halo outside the
cropped density support; the support weight itself is evaluated without that
clamp.  With $c_i(v)$ denoting the RGB query for vertex $v$, colour is blended
in floating point as

\begin{equation}
\bar{C}(v)=\frac{\sum_i w_i(v)c_i(v)}{\sum_i w_i(v)}.
\label{eq:colour-feather-fusion}
\end{equation}

A two-dimensional spatial index limits each window to nearby vertices, and
only one tri-plane is resident at a time.  Export requires positive colour
weight for every vertex, finite RGB values, unchanged vertex and face arrays,
and an exact colour round trip after reloading the PLY.

Algorithm~\ref{alg:vertex-colour-recovery} describes the second pass.  It uses
the geometry-fusion weights again, but accumulates RGB only for vertices whose
continuous coordinates lie inside a window's weighted support.

\begin{algorithm}[H]
\caption{Memory-bounded recovery of colour on a fixed global mesh}
\label{alg:vertex-colour-recovery}
\begin{algorithmic}[1]
\Require Mosaic $I$; fixed mesh $(V,F)$; window plan $S$; predictor $M$;
fusion metadata
\Ensure Per-vertex RGB array $C$ on unchanged $(V,F)$
\State Build a two-dimensional spatial index $B$ over the horizontal axes of $V$
\State Initialise RGB sum $R\gets 0$ and vertex weight sum $Z\gets 0$
\ForAll{windows $i\in S$}
  \State $T_i\gets\Call{EncodeWindow}{M,I,i}$ \Comment{one resident tri-plane}
  \State $(J,w)\gets\Call{SupportedVertices}{V,B,i,\text{fusion metadata}}$
  \ForAll{paired batches $(Q,w_Q)$ drawn from $(J,w)$}
    \State $q\gets\Call{ModelCoordinates}{V[Q],i}$
    \State $c\gets\Call{QueryColour}{M,T_i,q}$
    \State $R[Q]\gets R[Q]+w_Qc$
    \State $Z[Q]\gets Z[Q]+w_Q$
  \EndFor
  \State Release $T_i$
\EndFor
\State \textbf{require} $Z(v)>0$ and finite for every $v\in V$
\State $C\gets\Call{QuantiseRGB}{R/Z}$
\State \Call{ExportAndReloadVerify}{$V,F,C$}
\State \Return $C$
\end{algorithmic}
\end{algorithm}

\begin{figure}[H]
\centering
\resizebox{0.98\textwidth}{!}{%
\begin{tikzpicture}[
  node distance=7mm and 6mm,
  source/.style={draw=evidenceblue,rounded corners,align=center,
    text width=25mm,minimum height=10mm,fill=evidenceblue!8},
  stage/.style={draw=uclblue,rounded corners,align=center,
    text width=27mm,minimum height=10mm,fill=lightgrey},
  output/.style={draw=evidencegreen,rounded corners,align=center,
    text width=27mm,minimum height=10mm,fill=evidencegreen!8},
  arrow/.style={-{Latex[length=2mm]},thick,uclblue}
]
\node[source] (mosaic) {Rectangular zoom-20 mosaic};
\node[stage,right=of mosaic] (windows) {Overlapping source windows};
\node[stage,right=of windows] (patches) {Cropped density patches};
\node[stage,right=of patches] (fusion) {Fractional splat and feather};
\node[output,right=of fusion] (field) {Global density and one MC};
\draw[arrow] (mosaic) -- (windows);
\draw[arrow] (windows) -- (patches);
\draw[arrow] (patches) -- (fusion);
\draw[arrow] (fusion) -- (field);

\node[stage,below=11mm of windows] (triplane) {Recomputed tri-plane per window};
\node[stage,right=of triplane] (query) {Supported-vertex colour query};
\node[stage,right=of query] (blend) {Weighted floating-point RGB};
\node[output,right=of blend] (mesh) {Fixed geometry with vertex colour};
\draw[arrow] (windows) -- (triplane);
\draw[arrow] (triplane) -- (query);
\draw[arrow] (query) -- (blend);
\draw[arrow] (field) -- (mesh);
\draw[arrow] (blend) -- (mesh);
\end{tikzpicture}%
}
\caption{Two-pass large-image method.  The first pass blends density before a
single surface extraction; the second revisits each window to blend colour on
the unchanged global vertices.}
\label{fig:large-image-method-flow}
\end{figure}

\subsection{Independent-mesh stages and relation to Sat3DGen}

The evaluated sequence loads configuration, OSM, and DSM data; discovers OBJ
tiles; parses, crops, transforms, classifies, and pre-aligns each tile; removes
lower geometry; stitches vertex groups; applies DSM correction; and exports
paired DSM and no-DSM building branches.  The experiment isolates these
post-inference operations from upstream generation and optional fa\c{c}ade
processing.

Table~\ref{tab:sat3dgen-method-comparison} distinguishes the released
large-image wrapper, the density/colour extension, and the independent-mesh
geospatial path.

\begin{table}[H]
\centering
\caption{Relationship among the released Sat3DGen large-image wrapper, the
density/colour extension, and mesh-space geospatial integration.}
\label{tab:sat3dgen-method-comparison}
\footnotesize
\begin{tabularx}{\textwidth}{p{0.16\textwidth}XXX}
\toprule
Aspect & Released large-image wrapper & Density/colour extension &
Mesh-space and ChordAtlas path\\
\midrule
Learned representation & Released DINOv3-conditioned tri-plane density field &
Same checkpoint and model transform; no retraining & Released model or cached
OBJ enters explicit local coordinates\\
Overlap & Integer-aligned cropped fields and equal-weight averaging &
Fractional origins, bilinear placement, and raised-cosine weights & Independent
OBJ tiles are cropped, placed, and stitched\\
Surface and colour & One global Marching Cubes surface; large-image PLY has no
colour & One global surface followed by weighted per-vertex RGB recovery & RGB
OBJ assets, building extraction, and a separate ChordAtlas colour renderer\\
Geospatial context & Normalised scene coordinates & Rectangular mosaic and
explicit image-to-density scale & X-east/Y-up/Z-south frame, OSM footprints,
and EPSG:27700 DSM\\
\bottomrule
\end{tabularx}
\end{table}

Each stage was quantified using OBJ header counts and, where necessary,
streamed vertex comparison.  The following metrics were computed:

\begin{itemize}
  \item vertices and faces before crop, after crop, after bottom removal, after
  stitching, and after building extraction;
  \item geographic cluster count and separation from filename coordinates;
  \item number and distribution of per-vertex DSM $\Delta Y$ values;
  \item finite vertices, valid indices, zero-area and duplicate faces, boundary
  and non-manifold edges, connected components, and edge-incidence
  watertightness for a representative paired building; 
\end{itemize}

\subsection{Controlled tests of geometry rules}

Two small synthetic fixtures tested whether documentation and configuration
were enforced.  For bottom removal, one triangle had only one vertex within
\SI{0.5}{\metre} of the global minimum and a second triangle was fully elevated.
The stated all-three-vertices rule predicts that neither should be removed.  The
active function was invoked without modifying source code and its returned
face set recorded.

For stitching, two identically wound triangles had coincident X/Z coordinates
but were separated by \SI{10}{\metre} in Y.  They were assigned different tile
ranges and stitched with an X/Z radius of \SI{0.01}{\metre}.  The experiment
recorded vertex/face counts and retained attributes.  It directly tests whether
the configured maximum Y difference participates in matching.

\section{ChordAtlas integration protocols}

\subsection{Static workspace}

The static bridge combines a target bounding box, geospatial footprints,
existing mesh tiles, panoramas, and DSM prerequisites in a staged workspace.
Python records the manifest and local-frame footprint geometry; Java constructs
the ChordAtlas object graph and converts OBJ geometry into indexed MiniMesh
tiles.  A stored source offset preserves the original centroid during
conversion instead of implicitly recentering the input.

Validation checks required files, footprint geometry, MiniMesh structure, and
coordinate metadata.  Mesh coverage was independently intersected with the
target AOI. 

\subsection{On-demand selection and publication}

The on-demand path is embedded in the existing ChordAtlas GIS generator.  A
selected polygon is canonicalised by
ring start, winding, holes, duplicate removal, and six-decimal rounding before
generating a stable selection identifier.  Processing is serialised on one
worker, and interface updates return to the GUI event thread.

The Java service writes an atomic request file and launches the Python bridge.
The request distinguishes independent satellite-tile processing from the
large-image mode.  The tile path derives a geographically anchored allow-list
and selects fresh or cached corrected geometry.  The large-image path plans a
rectangular zoom-20 mosaic, executes or validates the density and colour
stages, and converts the fused mesh
once into the project local frame.  Both routes apply DSM processing and
associate geometry with requested footprints.

Each publishable building receives an exact source-mesh manifest, geometry
after ground processing, and metadata in a staging tree.  Geometry ownership
and filtering are applied per footprint rather than reusing one
selection-wide mesh under several building identities.  In large-image mode,
Python additionally requires RGB on every OBJ vertex and verifies that DSM
processing preserves colour indices.  Java displays this OBJ through a
dedicated vertex-colour renderer; the per-building MiniMesh is generated as a
separate representation used by the ChordAtlas workspace.  The staged Python
publication is installed atomically within its own boundary and marked
complete before Java conversion and renderer construction.

Figure~\ref{fig:publication-lifecycle} makes the transaction boundary explicit.

\begin{figure}[H]
\centering
\begin{tikzpicture}[
  node distance=7mm,
  state/.style={draw=uclblue,rounded corners,align=center,text width=38mm,
    minimum height=9mm,fill=lightgrey},
  good/.style={draw=evidencegreen,rounded corners,align=center,text width=38mm,
    minimum height=9mm,fill=evidencegreen!8},
  arrow/.style={-{Latex[length=2mm]},thick}
]
\node[state] (select) {GIS polygon and stable ID};
\node[state,right=of select] (python) {Python plan, inputs, generation, DSM, extraction};
\node[good,right=of python] (ready) {Python publication\\complete};
\node[state,below=of ready] (mini) {Java per-building MiniMesh conversion};
\node[good,left=of mini] (gui) {GUI building-generator creation};
\draw[arrow] (select) -- (python);
\draw[arrow] (python) -- (ready);
\draw[arrow] (ready) -- node[right,font=\scriptsize]{outside Python commit} (mini);
\draw[arrow] (mini) -- (gui);
\draw[dashed,evidenceorange,thick] ($(ready.south west)+(-1mm,-3mm)$) --
  ($(ready.south east)+(1mm,-3mm)$) node[midway,below=1mm,font=\scriptsize]
  {atomicity boundary};
\end{tikzpicture}
\caption{On-demand publication lifecycle.  Python completion precedes
the Java conversion required by ChordAtlas, so end-to-end publication is not a
single atomic transaction.}
\label{fig:publication-lifecycle}
\end{figure}

\subsection{Selection-scoped appearance inputs}

The appearance branch reuses the same geographic selection and local frame as
the building publication.  For roofs, the bridge selects the satellite tiles
already bound to the request, mosaics them in north-up orientation, clips them
to the building footprint, and publishes a reference image and style
mask.  Existing ready references can be added to an earlier building
publication without modifying its geometry.  ChordAtlas validates the
reference metadata before exposing it to the fa\c{c}ade and roof style
interface.

For street-level context, candidate viewpoints are sampled around the selected
OSM footprint and converted to WGS84.  Street View metadata are resolved and
deduplicated before acquisition.  The service presents one sample for approval,
downloads the remaining accepted batch, promotes the complete set, and writes
the panorama positions into ChordAtlas's local coordinate frame.  Figure~\ref{fig:appearance-input-flow}
summarises the two paths.

\begin{figure}[H]
\centering
\resizebox{0.98\textwidth}{!}{%
\begin{tikzpicture}[
  node distance=7mm and 6mm,
  source/.style={draw=evidenceblue,rounded corners,align=center,
    text width=27mm,minimum height=10mm,fill=evidenceblue!8},
  stage/.style={draw=uclblue,rounded corners,align=center,
    text width=29mm,minimum height=10mm,fill=lightgrey},
  output/.style={draw=evidencegreen,rounded corners,align=center,
    text width=30mm,minimum height=10mm,fill=evidencegreen!8},
  arrow/.style={-{Latex[length=2mm]},thick,uclblue}
]
\node[source] (sat) {Selection-bound\\satellite tiles};
\node[stage,right=of sat] (mosaic) {North-up mosaic\\and footprint mask};
\node[output,right=of mosaic] (roof) {Hashed roof\\reference};
\node[stage,right=of roof] (style) {ChordAtlas style\\input contract};
\draw[arrow] (sat) -- (mosaic);
\draw[arrow] (mosaic) -- (roof);
\draw[arrow] (roof) -- (style);

\node[source,below=11mm of sat] (footprint) {Selected OSM\\footprint};
\node[stage,right=of footprint] (viewpoints) {Viewpoint sampling\\and metadata deduplication};
\node[stage,right=of viewpoints] (approval) {Sample approval\\and batch acquisition};
\node[output,right=of approval] (panos) {Local-frame\\panorama index};
\draw[arrow] (footprint) -- (viewpoints);
\draw[arrow] (viewpoints) -- (approval);
\draw[arrow] (approval) -- (panos);
\end{tikzpicture}%
}
\caption{Selection-scoped appearance-input workflows.  The roof path prepares
a footprint-matched satellite reference for the style interface; the panorama
path acquires and registers street-level observations in the same local frame.}
\label{fig:appearance-input-flow}
\end{figure}

Selection logs and result manifests were classified by processing state:
request rejection, planning, acquisition, inference, corrected-scene reuse,
building publication, Java conversion, and GUI creation. 

\section{Matched representation comparison}

The directly matched comparison fixes OSM way~369245408 as the main-building
identity.  The large-image and independent-tile requests contain the same
target bounds and the same two footprint geometries.  The former begins from
one fused density field; the latter begins from nine independently extracted
meshes and adds OSM pre-alignment, lower-surface removal, stitching, and DSM
processing.  Building-level vertices, faces, relief, edge incidence, and
vertex-colour variation are compared under this shared spatial request.

Cross-system inspection is restricted to representation.  A retained
CityEngine 2026 export records one named initial shape, UV coordinates, and
UV-mapped materials.  The two
route-labelled ChordAtlas exports have envelopes consistent with the main
footprint; the standalone
ChordAtlas example has matched location.  All examples are
rendered with a common orthographic camera and metric scale within each figure.
The audit records polygon or triangle representation, vertex colour, texture
coordinates, and diffuse/normal/specular map organisation.  These properties
are not treated as fidelity scores because no surveyed geometry or registered
street-level appearance reference is available.

\section{Verification protocol}

\subsection{Executable tests}

Component tests cover configuration, coordinate frames, footprints, DSM
prerequisites, OBJ inspection, tile planning, image integrity, manifests,
cache freshness, publication, rollback, MiniMesh conversion, OBJ slicing, and
vertex-colour parsing.  Separate large-image regression cases exercise
fractional placement, non-zero fusion coverage, seam continuity, colour-query
coordinates, weighted RGB, and PLY round trips.  The bridge suite also covers
large-image selection, roof-reference and panorama planning, promotion, and
coordinate transfer.  Python, algorithm, and Java tests are reported
separately because they exercise different runtime and data-transfer
boundaries.  Detailed results are provided in
Appendix~\ref{app:evidence-record}.

\subsection{Validity review}

Construct validity was assessed by matching each measure to the property it
was intended to represent.  The review distinguishes geographic extent from
triangle coverage, file availability from semantic correctness, topological
validity from geometric accuracy, and producer-side publication from
consumer-side conversion.  External validity is examined through the spatial
and dataset coverage of the case study.

%%%%%%%%%%%%%%%%%%%%%%%%%%%%%%%%%%%%%%%%%%%%%%%%%%%%%%%%%%%%%%%%%%%%%%%%%%%%%%%%
\chapter{Results and Analysis} \label{Chap4}\label{chap:results}

\section{Evaluation scope}

Results are organised around the evaluation questions: data integrity,
large-image density and colour inference, independent-mesh behaviour,
coordinate control, static and on-demand integration, matched representation,
and component verification.  Quantitative findings are tied to the data population or
processing stage from which they were measured.

\section{Prototype feasibility and empirical focus}

A retained London GLB and corresponding server record establish compatibility
with the released Sat3DGen inference interface.  The Unity/Ubiq client also
compiled with the intended acquisition, request, retrieval, and scene-loading
stages, informing the design of the downstream interface.  The principal
experiments comprise a fresh contiguous London large-image run, a matched
overlap-fusion comparison, global vertex-colour recovery, an independent-OBJ
mesh-processing study, and ChordAtlas publication.  Retained ChordAtlas
appearance exports provide representation-level evidence, while the earlier
format-conversion prototypes contribute architectural context.

\section{Data construction and integrity}\label{sec:data-results}

The London builder is assessed for dataset readiness and downstream geospatial
reuse.  The results address coverage, independence, label geometry, and
auxiliary-channel integrity rather than model-training performance.

\subsection{Inventory and physical extent}

Panorama acquisition covered the majority of target locations, while the satellite, DSM, OSM, and grid channels cover only a subset; sampling density and overlap also differ between branches.

Projected tile unions expose the difference between file count and physical
coverage.  The 127-tile set covers approximately
\SI{0.45}{\kilo\metre\squared}; the repaired grid covers
\SI{0.96}{\kilo\metre\squared}.  The 50 and 10 per cent overlap layouts cover
approximately \SI{0.93}{\kilo\metre\squared} and
\SI{0.89}{\kilo\metre\squared}, respectively.  Despite containing 306 tiles,
the latter supports 2,331/2,333 queries because its stride and boundary handling
omit two locations.  Tile count is therefore an insufficient coverage measure.

\begin{figure}[H]
\centering
\includegraphics[width=\textwidth]{data_builder_coverage_channels.pdf}
\caption{Estimated grid-footprint union, matched-query coverage, and stored
channel availability.  Areas use a local-scale Web-Mercator estimate.}
\label{fig:data-builder-coverage-channels}
\end{figure}

\subsection{DSM channel}

The four Environment Agency inputs are adjacent 5,000-by-5,000 float rasters at
\SI{1}{\metre} resolution in EPSG:27700.  The product is the Composite 2022 DSM
\cite{environmentAgencyDSM2022}; it measures visible surfaces, not bare earth.
The source rasters cover all 2,343 query points and all 272 repaired
grid centres.

Figure~\ref{fig:data-builder-dsm-views} visualises the metric source mosaic
before per-patch percentile normalisation.  Mosaicking the four rasters across
the approximately 1.00-by-0.94-kilometre footprint provides complete coverage
at the displayed sampling grid.

\begin{figure}[H]
\centering
\includegraphics[width=\textwidth]{data_builder_dsm_views.pdf}
\caption{North-up and oblique views of the Environment Agency 2022
\SI{1}{\metre} first-return DSM mosaic over the London study area.  The plan
view preserves equal east--north scales and NoData; the oblique view uses
threefold vertical exaggeration.  Colour shows elevation relative to Ordnance
Datum Newlyn.  Source: Environment Agency, LIDAR Composite DSM 2022; Open
Government Licence v3.0
\cite{environmentAgencyDSM2022,nationalArchivesOGL3}.}
\label{fig:data-builder-dsm-views}
\end{figure}

Independent 5th--95th-percentile scaling makes the patches suitable as relative
surface-height channels but removes cross-patch comparability in absolute
elevation.  A legacy DSM derivative without traceable production metadata was
therefore excluded from the evaluated channel.

\subsection{OSM semantic channel}

The Overpass query retrieves support nodes recursively, and the extractor emits
all returned nodes as point features even when they have no category tags.  Of
67,281 category records, 57,623 are such support-node records (85.65 per cent).
Multipolygon fragments are also treated as
independent outer rings, and inner-ring assignment can fill intended holes.
Semantic validity therefore requires geometry- and tag-aware validation rather
than file-count checks.
The active mesh loader adds 19,682 ordinary-building and 9,237 height-building
records to the same index
\cite{openstreetmapData}.

\section{Large-image density and colour inference}
\label{sec:large-image-results}

\subsection{Contiguous London configuration}

The contiguous run covered approximately \SI{285.5}{\metre} by
\SI{475.8}{\metre} from a 3,072-by-5,120-pixel zoom-20 mosaic.  Each
640-pixel source window was passed through the unchanged 256-pixel Sat3DGen
image transform.  A 160-pixel stride produced a 17-by-29 schedule, or 493
windows, with 75 per cent nominal overlap.  The mesh-resolution setting was
192 voxels; it did not change the network image size.  After fractional
placement, the fused density volume measured $1232\times740\times173$ and had
no zero-weight cells.

One global Marching Cubes extraction produced 3,122,922 vertices and 6,241,862
faces.  Density inference and fusion required 129.0 seconds and reached
\SI{6.62}{\gibi\byte} peak allocated GPU memory.

\subsection{Window-boundary comparison}

The hard-box and fractional-feather runs used the same source image,
checkpoint, windows, overlap, voxel resolution, and isovalue.  The baseline is
described as release-style because it retains the released wrapper's hard
boundary crop and equal-weight overlap rule, while using the matched project
settings rather than the release defaults.  Figure~\ref{fig:large-image-feather}
shows the resulting top surfaces.

\begin{figure}[H]
\centering
\includegraphics[width=\textwidth]{big_image_feather_comparison.pdf}
\caption[Hard-box and fractional-feather large-image fusion.]{Matched
large-image fusion comparison.  Hard-box averaging produces repeated
window-aligned height steps; fractional placement and raised-cosine feathering
replace the abrupt contributor changes with smooth overlap transitions.  The
false colour and shading encode density-derived top-surface height, not
satellite texture or appearance accuracy.}
\label{fig:large-image-feather}
\end{figure}

At the former hard-crop contribution boundaries, mean absolute density jumps
fell from 0.710 to 0.286 along rows and from 0.963 to 0.491 along columns
(Table~\ref{tab:large-image-seams}).  Mean top-surface jumps fell from 2.596 to
0.659 voxels and from 3.065 to 1.214 voxels, respectively.  The corresponding
reductions are 49--60 per cent for density and 60--75 per cent for the derived
top surface.  Geometry runtime increased by 8.6 per cent, while peak allocated
GPU memory was unchanged at the reported precision.  The comparison therefore
supports a reduction in boundary discontinuity.

\begin{table}[H]
\centering
\caption{Matched hard-box and fractional-feather fusion results.}
\label{tab:large-image-seams}
\small
\begin{tabularx}{\textwidth}{Xrrr}
\toprule
Measure & Hard-box & Fractional feather & Change\\
\midrule
Mean density jump, row boundary & 0.710 & 0.286 & $-59.8\%$\\
Mean density jump, column boundary & 0.963 & 0.491 & $-49.0\%$\\
Mean top-surface jump, row (voxels) & 2.596 & 0.659 & $-74.6\%$\\
Mean top-surface jump, column (voxels) & 3.065 & 1.214 & $-60.4\%$\\
Geometry stage (s) & 118.7 & 129.0 & $+8.6\%$\\
Peak allocated GPU memory (GiB) & 6.62 & 6.62 & unchanged\\
\bottomrule
\end{tabularx}
\end{table}

\subsection{Global vertex-colour recovery}

The second pass accumulated 42.6 million positive window contributions over
the fixed global surface.  Every vertex received support from at least one
window, with no uncovered or all-black vertices.  Reloading the coloured PLY
preserved every vertex position and face index and reproduced the exported RGB
values.  The colour stage required 97.3 seconds, of which 80.9 seconds were
model queries, and reached \SI{1.86}{\gibi\byte} peak allocated GPU memory.
The exported PLY displayed its blended vertex colours in MeshLab.  These checks
establish complete and index-consistent appearance transfer.

\subsection{Current myProject large-image publication}

The large-image path was also exercised through the ChordAtlas bridge on a
1,536-by-1,536-pixel, 49-window selection.  Its retained job records a
340,825-vertex, 678,726-face density surface followed by complete colour
support.  The geometry stage required 17.0 seconds and the colour stage 23.2
seconds.  The final colour-enabled publication reused the validated geometry
stage from the same retained job before running and validating the second pass.

After one local-coordinate conversion, the context crop contained 66,666
vertices and 132,440 faces.  DSM processing covered every retained vertex and
preserved its RGB index while changing vertical position by at most
\SI{30.7}{\metre}.  Footprint extraction published two non-exclusive coarse
building assets.  Their geometric union, represented by the top-level
selection crop, contains 17,571 vertices and 34,383 faces; the two building
files share 266 vertex positions and 452 triangles.  Both OBJ files contain
RGB on every vertex and pass the Java colour parser; per-building MiniMesh
trees also exist.  In ChordAtlas, the RGB OBJ is displayed by a dedicated
vertex-colour renderer, while MiniMesh remains a separate geometric workspace
representation.

\begin{figure}[H]
\centering
\includegraphics[width=\textwidth]{big_image_coloured_publication.pdf}
\caption{Oblique and plan views of the current large-image building
publication after coordinate conversion, DSM processing, and footprint
extraction.  Face colour is the mean of the predicted per-vertex RGB values;
the figure demonstrates colour retention through publication, not comparison
with a reference texture.}
\label{fig:large-image-coloured-publication}
\end{figure}

\section{Mesh-space post-inference experiment}\label{sec:sat-results}

\subsection{Evaluated processing scope}

The experiment processes 15 coordinate-paired OBJ tiles through OSM/DSM
loading, crop and placement, lower-surface removal, stitching, DSM correction,
and building extraction.  The associated cache records describe geometry but
do not bind a model checkpoint and configuration, so the measurements below
characterise the mesh-processing stages.

Qian et al. report FID 19.2 and DSM RMSE \SI{5.20}{\metre} on VIGOR-OOD
\cite{qian2026sat3dgen}.  Those metrics assess the published model protocol;
this experiment measures downstream operations on retained mesh tiles and uses
the external DSM as a processing input.  Unlike the density-domain experiment
in Section~\ref{sec:large-image-results}, it combines independently extracted
meshes and therefore evaluates a complementary assembly strategy.

\subsection{Geographic composition of the mesh batch}

Recursive discovery selected two disconnected clusters: nine eastern tiles in
a 3-by-3 arrangement and six western tiles in a 2-by-3 arrangement.  Their
centres are approximately \SI{666}{\metre} apart, more than eleven nominal tile widths.
Figure~\ref{fig:tile-clusters} plots every centre and nominal footprint.

\begin{figure}[H]
\centering
\includegraphics[width=0.94\textwidth]{tile_clusters.pdf}
\caption{Geographic composition of the 15 OBJ meshes selected by the merge
experiment.  Filename coordinates identify two disconnected sites.}
\label{fig:tile-clusters}
\end{figure}

The merged output consequently combines two sites rather than one contiguous
London scene.  This geographic inconsistency originates in cache selection,
not the upstream Sat3DGen generator \cite{qian2026sat3dgen}.

\subsection{Geometry retained through processing}

Table~\ref{tab:stage-counts} gives the exact geometry counts.  From raw inputs
to the stitched scene, 1,724,759/3,640,214
vertices (47.38 per cent) and 3,184,234/7,281,608 faces (43.73 per cent) were
removed or consolidated.  This large, intended reduction makes the crop,
bottom-removal, and stitching predicates central to output validity.

\begin{table}[H]
\centering
\caption{Geometry counts by post-processing stage.}
\label{tab:stage-counts}
\small
\begin{tabular}{lrrr}
\toprule
Stage & Vertices & Faces & Change from previous stage\\
\midrule
Raw 15 OBJ inputs & 3,640,214 & 7,281,608 & ---\\
Cropped & 2,828,613 & 5,615,039 & $-811{,}601$ V; $-1{,}666{,}569$ F\\
Bottom removed & 2,252,373 & 4,474,289 & $-576{,}240$ V; $-1{,}140{,}750$ F\\
Stitched & 1,915,455 & 4,097,374 & $-336{,}918$ V; $-376{,}915$ F\\
DSM scene & 1,915,455 & 4,097,374 & topology unchanged\\
DSM building extraction & 1,001,507 & 2,116,136 & branch result\\
No-DSM building extraction & 907,370 & 1,911,272 & branch result\\
\bottomrule
\end{tabular}
\end{table}

Figure~\ref{fig:tile09-bottom-removal} shows the saved crop-stage and
bottom-removal outputs for one representative tile.  The horizontal footprint
is unchanged, while the operation removes a near-horizontal generated support
sheet spanning almost the full tile.  Every tile lost exactly 38,416 vertices
and 76,050 faces at this stage, consistent with removal of the same standard
lower grid.

\begin{figure}[H]
\centering
\includegraphics[width=0.98\textwidth]{tile09_bottom_removal.pdf}
\caption[Bottom removal in representative tile 09.]{Representative tile~09 at
two consecutive saved stages.  Panel (a) is the cropped, placed, and
OSM-pre-aligned mesh before bottom removal, with the removed lower sheet
highlighted in orange; panel (b) is the retained mesh immediately afterwards.
Both use the same underside camera and metric frame, and stored vertex colour
is omitted to expose the geometric change.}
\label{fig:tile09-bottom-removal}
\end{figure}

The controlled bottom-removal fixture also showed that the active predicate
removes a triangle when only one vertex is close to the lowest surface,
although the documented rule requires all three vertices to satisfy the
threshold.  In the stitching fixture, two triangles separated by
\SI{10}{\metre} in Y collapsed from six vertices to three because matching used
X/Z coincidence without the configured vertical-separation test.  In the
north-western two-by-three subset used below, the six pre-stitch tiles contain
819,191 vertices and 1,626,029 faces; the saved post-stitch subset contains
696,199 vertices and 1,485,998 faces with the same horizontal extent.  The
sub-metre consolidation is quantified here because it is not distinguishable
in a scene-scale silhouette.

\subsection{DSM displacement is scene-wide}

DSM correction changed all 1,915,455 vertices in the saved scene and every
change was positive.  The median displacement is \SI{36.08}{\metre}, the mean
is \SI{42.24}{\metre}, and values range from \SI{15.81}{\metre} to
\SI{62.63}{\metre}.  The scene's vertical range changes from
\SIrange{-26.77}{18.76}{\metre} to \SIrange{1.86}{67.94}{\metre}
(Figure~\ref{fig:dsm-displacement}).

\begin{figure}[H]
\centering
\includegraphics[width=\textwidth]{dsm_displacement.pdf}
\caption{Per-vertex vertical displacement introduced by DSM correction.  The
distribution characterises implementation behaviour rather than surveyed
height error.}
\label{fig:dsm-displacement}
\end{figure}


\subsection{North-western stitched scene and representative topology}

Figure~\ref{fig:dsm-scene-nw-2x3} isolates the north-western two-by-three
cluster in the stitched scene before and after the saved DSM stage.  The two
OBJ files preserve vertex order, horizontal coordinates, RGB attributes, and
face indices within this subset, so the stage changes only Y.  The subset
contains 696,199 vertices and 1,485,998 faces.  All paired vertices move
upward; median displacement is \SI{38.58}{\metre}, with a 5th--95th percentile
range of \SIrange{32.96}{60.49}{\metre}.

\begin{figure}[H]
\centering
\includegraphics[width=0.98\textwidth]{dsm_scene_nw_2x3_comparison.pdf}
\caption[DSM-stage vertical transformation in the north-western two-by-three
cluster.]{DSM-stage transformation in the north-western two-by-three stitched
scene cluster.  Panel (a) is the post-stitch output before DSM processing;
panel (b) is the same mesh after DSM height correction.  Both use uniform
materials, a shared camera, and a common metric frame.  Panel (c) reports
paired $\Delta Y$ at the highest pre-correction vertex in each occupied
1-metre X/Z cell.  Stored RGB is omitted to expose the vertical change; the
values describe implemented behaviour rather than surveyed elevation error.}
\label{fig:dsm-scene-nw-2x3}
\end{figure}

A separate representative paired building remains a single connected
component with finite vertices, valid indices, and no detected zero-area,
duplicate, or non-manifold faces. Its mean paired vertical displacement is approximately
\SI{58.0}{\metre}.  This building-level case establishes the presence of open
boundaries; the population analysis in Section~\ref{sec:ondemand-results}
evaluates their prevalence.

\section{Static ChordAtlas integration}\label{sec:static-results}

The static path converted a high-resolution source mesh into a structurally
valid MiniMesh workspace with the required indices and generator structure.
The mesh's rectangular X/Z extent overlaps approximately
\SI{0.0096}{\kilo\metre\squared} of a \SI{1.00}{\kilo\metre\squared} target,
or 0.96 per
cent.  This is an axis-aligned extent proxy, not surface area, so it can include
empty regions.  Footprint-associated assets were exported, with overlapping
OSM filter categories reported separately.

The evaluated result demonstrates

\begin{center}
cached meshes $\rightarrow$ coordinate/footprint processing $\rightarrow$
MiniMesh conversion $\rightarrow$ structurally valid workspace,
\end{center}

using retained inputs.  It establishes coordinate, footprint, format, and
workspace integration over the available geographic extent.

\section{On-demand ChordAtlas integration}\label{sec:ondemand-results}

\subsection{Independent-tile publication population and asset quality}

The tile-based on-demand set contains ten publication records with
building-level metrics and one earlier record without them.
Table~\ref{tab:on-demand-cases} aggregates these records by corrected-scene
decision.

\begin{table}[H]
\centering
\caption{Independent-tile publication records by corrected-scene decision.}
\label{tab:on-demand-cases}
\small
\begin{tabularx}{\textwidth}{XrrrX}
\toprule
Record category & Publications & Tile references & Buildings & Publication evidence\\
\midrule
Fresh correction & 5 & 90 & 11 & Modular building assets published\\
Reused correction & 5 & 91 & 8 & Modular building assets published\\
Earlier schema & 1 & 9 & --- & Publication boundary reached; no current metrics\\
\bottomrule
\end{tabularx}
\end{table}

Across the ten metric-bearing independent-tile records, 181 tile references produced 19
building assets: eleven following fresh correction and eight following cache
reuse.  Figure~\ref{fig:on-demand-publications-update} distinguishes selection
size from the number of modular outputs.

\begin{figure}[H]
\centering
\includegraphics[width=\textwidth]{on_demand_publications_update.pdf}
\caption[On-demand publication records and corrected-scene reuse.]{On-demand publication records, tile counts, building counts, and
corrected-scene reuse decisions.  The Python publication boundary precedes
Java conversion.}
\label{fig:on-demand-publications-update}
\end{figure}

Topology is the principal quality constraint in this published population.

\begin{figure}[H]
\centering
\includegraphics[width=\textwidth]{on_demand_building_topology_update.pdf}
\caption{Boundary-edge metrics for all 19 independent-tile building publications and
non-manifold-edge metrics for the ten v2 publications that record this field.}
\label{fig:on-demand-topology-update}
\end{figure}

RGB attributes are preserved throughout the independent-tile publication set:
every vertex in all 19 cropped OBJ assets carries colour data.  The
representative output in Figure~\ref{fig:on-demand-vertex-colour} shows this
appearance together with the open-edge topology identified above.

All 19 independent-tile buildings also contain MiniMesh indices.  A real-asset Java check parsed
vertex colours and constructed the generator layer successfully, confirming
appearance transfer through OBJ parsing and generator construction.

\begin{figure}[H]
\centering
\includegraphics[width=\textwidth]{on_demand_vertex_colour_building.pdf}
\caption{Oblique and plan views of a representative building-level
vertex-colour OBJ.  Colour visualises the published appearance attributes.}
\label{fig:on-demand-vertex-colour}
\end{figure}


\subsection{Spatially matched appearance references}

The independent-tile bridge extends building publication with two forms of
appearance input.  Four building assets have ready satellite roof references generated
from the tiles bound to their selections.  Each reference covers its source
footprint completely, records a footprint-derived style mask, and leaves the
published geometry unchanged.  This creates a spatially matched roof-appearance
input for the ChordAtlas style interface.

A selected-block panorama run also completed its planning, sample approval,
batch acquisition, and promotion stages.  Eight unique Street View panoramas
were acquired and indexed in the ChordAtlas workspace using the project's
local coordinate frame.  The result connects a chosen OSM footprint to a
deduplicated, selection-specific set of street-level observations.

Together, these results extend the bridge from geometry delivery to registered
appearance inputs.  For this reference-registration experiment, the
roof-reference style contract and panorama services pass their component
checks; the evaluated outcome is reference preparation and spatial
registration rather than a texture-quality measure.

\subsection{Failure containment}

Input validation rejected a duplicate footprint identity and an incomplete
imagery set before inference.  Both cases reached explicit terminal states,
demonstrating \gls{failclosed} behaviour at the request and input-validation
boundaries.  Append-only selection logs preserve these transitions.


\section{Matched London building comparison}
\label{sec:matched-building-comparison}

\subsection{Directly matched Sat3DGen routes}

The large-image and independent-tile manifests request the same target bounds
and the same two OSM footprints.  Figure~\ref{fig:matched-sat3dgen-building}
compares the main building, OSM way~369245408, with identical cameras and
metric display scale.  Both outputs store per-vertex RGB produced within their
respective Sat3DGen routes, and both
were processed against the same one-metre DSM with complete recorded vertex
coverage.  Table~\ref{tab:matched-building-metrics} reports corresponding
geometry, topology, relief, and colour-variation descriptors.

\begin{figure}[!htbp]
\centering
\includegraphics[width=\textwidth]{matched_sat3dgen_building_comparison.pdf}
\caption[Matched Sat3DGen building comparison.]{Directly matched building
comparison for OSM way~369245408.  Panels (a,c) show the fractional-feather
large-image publication and panels (b,d) the nine-mesh independent-tile
publication.  The latter retains additional sheets, spikes, and detached
fragments in this sample.  Colours are per-vertex RGB retained from the two
Sat3DGen routes, not satellite pixels projected directly onto the mesh.}
\label{fig:matched-sat3dgen-building}
\end{figure}

\begin{table}[H]
\centering
\caption{Descriptive metrics for the directly matched main-building outputs.}
\label{tab:matched-building-metrics}
\small
\begin{tabularx}{\textwidth}{p{0.18\textwidth}p{0.18\textwidth}p{0.10\textwidth}p{0.18\textwidth}p{0.12\textwidth}X}
\toprule
Route & Vertices / faces & Relief (m) & Boundary / non-manifold & Luminance s.d. & Mean edge RGB $\Delta$\\
\midrule
Large-image & 16,644 / 32,597 & 28.12 & 681 / 0 & 0.154 & 0.053\\
Independent tiles & 21,555 / 42,524 & 38.07 & 1,074 / 384 & 0.213 & 0.082\\
\bottomrule
\end{tabularx}
\end{table}

The large-image sample has a more compact surface and fewer recorded boundary
and non-manifold edges.  The independent-tile sample contains more geometry
and stronger local RGB variation.  The latter observation is consistent with
visibly higher-frequency colour detail, but it is not evidence of more accurate
colour.  Similarly, complete DSM coverage shows that both outputs were
vertically conditioned.

\subsection{Procedural appearance representations}

Figure~\ref{fig:procedural-representation-examples} places the CityEngine export
and retained ChordAtlas artefacts under a common presentation camera and metric
scale.  These are illustrative representation examples rather than a
common reconstruction benchmark.  The route-labelled ChordAtlas envelopes are
geometrically consistent with the main footprint, but their OBJ files do not
embed its OSM identifier.  The standalone ChordAtlas model has neither a
matched location nor a retained run manifest.

\begin{figure}[!htbp]
\centering
\includegraphics[width=\textwidth]{procedural_representation_examples.pdf}
\caption[Procedural model representation examples.]{Illustrative procedural
representation examples, centred without rescaling.  The CityEngine export is
a compact UV-mapped envelope with repeated facade atlases.  The two
route-labelled ChordAtlas exports use generated UV materials rather than
vertex-RGB encoding and preserve the broad building envelope.  The standalone
ChordAtlas asset is included only as a stored
contextual reference.  Different or undocumented inputs prevent a fidelity
ranking across panels.}
\label{fig:procedural-representation-examples}
\end{figure}

The CityEngine 2026 log records one initial shape exported as 120 vertex
records, 49 polygon records, 84 texture coordinates, and three materials.  Its
facades are regular, and the visible doors and windows reside in repeated
texture atlases rather than separate door or window geometry.  This supports a
conclusion about export compactness and material organisation, not about
site-specific texture accuracy \cite{esriCityEngineIntro2026,
esriCityEngineSDK2026}.

The two ChordAtlas exports contain no vertex RGB.  Instead, the large-image
route stores 27 and the independent-tile route 47 complete diffuse, normal, and
specular map sets.  Window-labelled materials and the aligned map triplets
provide a multi-channel material representation rather than a single RGB value
at each vertex.  Figure~\ref{fig:appearance-map-encoding} illustrates that
distinction.  The tile-derived export also contains one material named
\texttt{error} on ten face records.

\begin{figure}[!htbp]
\centering
\includegraphics[width=\textwidth]{appearance_map_encoding_examples.pdf}
\caption[Facade appearance-map encodings.]{Selected facade material encodings.
The CityEngine facade material uses a repeated diffuse atlas; no normal or
specular map is stored for that facade material.  The retained ChordAtlas
procedural appearance exports store aligned diffuse, normal, and specular maps
for window-labelled regions \cite{kelly2018frankengan,kellyChordAtlas2024}.  These
maps can depict frames, panes, and material response.}
\label{fig:appearance-map-encoding}
\end{figure}

The retained route-labelled artefacts place satellite-route vertex colour and
structured procedural appearance maps side by side, although the OBJ files do
not provenance-lock a direct conversion between them.
They do not show that ChordAtlas preserves the Sat3DGen triangulation,
reconstructs the true facade texture, or produces a more accurate building.

\FloatBarrier

\section{Results synthesis}

All component tests passed within their specified scope, supporting the
coordinate, cache, geometry-processing, and publication contracts across
Python and Java.  Appendix~\ref{app:evidence-record} provides the detailed test
categories and results.  Tables~\ref{tab:results-synthesis} and
\ref{tab:results-synthesis-reliability} consolidate the principal findings by
evaluation dimension.

\begin{table}[H]
\centering
\caption{Synthesis of principal findings and their evidential scope.}
\label{tab:results-synthesis}
\small
\begin{tabularx}{\textwidth}{X p{0.20\textwidth}X}
\toprule
Evaluation dimension & Evidence basis & Finding and interpretation\\
\midrule
Static ChordAtlas conversion & Structural inspection & A cached high-resolution mesh produced a valid MiniMesh workspace.\\
London data integrity & Hash and spatial analysis & Exact panorama duplication, split leakage, candidate-offset errors, and OSM support-node contamination require correction before learning use.\\
Large-image reconstruction & Matched run and round-trip inspection & Feathering reduced window-boundary discontinuities; every global-mesh vertex received colour without changing geometry, and the current ChordAtlas path retained RGB through publication.\\
Matched building representation & Shared OSM request, OBJ/MTL audit, and common-camera rendering & The large-image sample is more compact, while the tile sample has stronger colour variation and more fragments; CityEngine and ChordAtlas differences support representation-level, not accuracy, conclusions.\\
Building and appearance publication & Execution and component tests & Nineteen independent-tile assets and two current large-image assets were published; the paths retain RGB and provide either registered appearance references or the second-pass colour renderer.\\
\bottomrule
\end{tabularx}
\end{table}

\begin{table}[H]
\centering
\caption{Synthesis of coordinate and geometry-reliability findings.}
\label{tab:results-synthesis-reliability}
\small
\begin{tabularx}{\textwidth}{X p{0.20\textwidth}X}
\toprule
Evaluation dimension & Evidence basis & Finding and interpretation\\
\midrule
Coordinate control & Calculation and projection comparison & The local transform is reversible, with up to approximately \SI{19.1}{\metre} disagreement relative to EPSG:27700 displacement.\\
Stitching rule & Controlled regression & A \SI{10}{\metre}-separated fixture collapsed from six to three vertices because vertical separation was omitted.\\
DSM correction rule & Artefact comparison & Every scene vertex moved positively.\\
Building topology & Population inspection & All 19 independent-tile publications and both current large-image publications retain boundary edges and none is watertight; nine of the ten v2 independent-tile assets reporting non-manifold incidence contain such edges.\\
\bottomrule
\end{tabularx}
\end{table}

%%%%%%%%%%%%%%%%%%%%%%%%%%%%%%%%%%%%%%%%%%%%%%%%%%%%%%%%%%%%%%%%%%%%%%%%%%%%%%%%
\chapter{Discussion} \label{Chap5}\label{chap:discussion}

\section{Interpretation of the research scope}

This dissertation examines how overlapping satellite-conditioned predictions
can form a continuous coloured regional surface and then become editable urban
assets.  It contributes at two interfaces: density-domain fusion and colour
recovery around the released Sat3DGen model, followed by geospatial association
and publication through OSM, DSM, and ChordAtlas.  This addresses the gap
between plausible window-level generation and spatially identified, reusable
urban content.

\section{Research question: extent and reliability boundary}

The study establishes a working large-image reconstruction and
coordinate-explicit integration at building-asset scale.  Overlapping density
windows are fused into one surface, per-vertex appearance is recovered in a
second pass, and the resulting mesh can be transformed, associated with OSM
footprints, processed against DSM context, and converted for ChordAtlas.
Stable selection identities, manifests, cache checks, and validation at the
Python--Java boundary make these transformations inspectable rather than
implicit.

Reliable regional operation depends on four properties highlighted by the
results.  Overlapping predictions need continuous placement and smooth
weighting; spatial predicates must prevent unrelated or
vertically separated surfaces from being combined; geometry processing must
preserve topology and apply external height information selectively; and
publication must be acknowledged by the consuming application.  These
requirements connect large-area reconstruction quality to asset-level system
reliability.

\section{Contribution of the ChordAtlas integration}

The contribution is defined by the new data and execution contracts introduced
at the ChordAtlas boundary.  The integration extends the existing GIS workflow
in the following ways:

\begin{itemize}
  \item a coordinate contract that preserves geographic selection, a local
  origin, and building identity across mesh conversion;
  \item a selectable large-image source mode that binds the mosaic plan,
  density fusion, colour recovery, coordinate conversion, and footprint crop;
  \item manifest-driven planning and cache validation that bind a request to
  its required tiles and processing stages;
  \item OSM- and DSM-aware processing that associates generated surfaces with
  footprints and publishes separately editable building assets;
  \item selection-scoped roof references and panoramas, exercised for the
  independent-tile selections, that connect appearance inputs to the same
  footprint identity and local coordinate frame;
  \item staged Python--Java publication with explicit failure states,
  background execution, and a dedicated per-vertex-colour renderer; and
  \item executable checks for coordinate conversion, mesh prerequisites,
  selection manifests, cache state, and publication contracts.
\end{itemize}

Together, these changes turn a fused or independently assembled surface into
geographically identified ChordAtlas assets whose processing history and
failure state can be examined.  Sat3DGen provides learned geometry and
appearance, while ChordAtlas provides asset structure and procedural editing.
The bridge preserves the coloured OBJ for display and builds MiniMesh as a
separate workspace representation; successful colour transfer should not be
conflated with colour being embedded in the MiniMesh tiles themselves.

For the independent-tile selections, the roof-reference and panorama results
are important because they associate
appearance observations with the same building identity used by the geometry
pipeline.  They provide registered inputs for ChordAtlas style generation.
The separately retained appearance exports establish the availability of
generated appearance artefacts, while reference registration, material
completeness, and texture fidelity
remain distinct evaluation questions.

\section{Relationship among the evaluated systems}

The project contributes at two interfaces to Sat3DGen.  First, it extends the
released large-image wrapper while retaining the pretrained DINOv3/tri-plane
representation and one global density-domain surface extraction
\cite{qian2026sat3dgen,qianSat3DGenCode2026}.  Fractional placement and smooth
overlap weighting address quantisation and hard contributor changes, while a
second pass restores the colour capability available in single-image
inference.  The matched hard-box experiment isolates the effect of this fusion
rule on boundary continuity and resource use.

Secondly, the mesh-space path introduces geographic placement, OSM/DSM
constraints, building extraction, and ChordAtlas publication.  This path
remains relevant where independently extracted or cached meshes are the
available interface.  The two experiments therefore compare representation
stages rather than treating mesh stitching as the only regional strategy.
Independent reference geometry and street-level imagery are still required to
evaluate metric and perceptual accuracy.

The retained London examples clarify that the systems address different
subproblems.  CityEngine generates controllable polygonal content from GIS
shapes, attributes, assets, and CGA rules and supports scripted procedural
runtimes \cite{esriCityEngineIntro2026,esriCityEngineGIS2026,
esriCityEngineSDK2026}.  Its retained export is a compact UV-mapped envelope
whose openings are represented primarily in repeated facade atlases.  Sat3DGen
instead conditions dense geometry and vertex appearance on overhead imagery.
The retained ChordAtlas exports organise surfaces into UV material groups and
appearance maps.  The bridge connects these representations through stable
identity, coordinates, and editable asset boundaries.

Table~\ref{tab:paradigm-comparison} summarises the distinct evidence and role of
each system in this study.

\begin{table}[H]
\centering
\caption{Evaluated systems, output representations, and their roles in the
study.}
\label{tab:paradigm-comparison}
\footnotesize
\begin{tabularx}{\textwidth}{>{\raggedright\arraybackslash}p{0.16\textwidth}
>{\raggedright\arraybackslash}X>{\raggedright\arraybackslash}X
>{\raggedright\arraybackslash}X}
\toprule
System & Input and control & Output representation & Evidence and role\\
\midrule
CityEngine / CGA & Platform support for GIS shapes, CGA rules, attributes, and
assets; one initial shape in the retained case & Compact UV-mapped polygon
export & Illustrative export for compactness and material organisation\\
Sat3DGen large-image & Rectangular mosaic, overlap schedule, and fusion weights
& One density-fused surface with predicted vertex RGB & Directly matched OSM
case; evaluated for continuity, colour support, and topology\\
Sat3DGen independent tiles & Nine image-conditioned meshes plus OSM/DSM
operations & Stitched per-footprint assets with stored vertex RGB & Directly
matched OSM case; evaluated for mesh-space assembly behaviour\\
ChordAtlas procedural appearance & Structured procedural surfaces and style
controls & Semantic UV groups with diffuse, normal, and specular maps &
Retained procedural appearance representation; not reconstruction truth or topology
preservation\\
Hybrid bridge & Geographic selection, manifests, OSM/DSM context, and
publication controls & Identified coloured OBJ, MiniMesh workspace, and
registered appearance inputs & Integration contribution across learned and
procedural representations\\
\bottomrule
\end{tabularx}
\end{table}

This is a partially matched representation case rather than a common reconstruction
benchmark.  The two Sat3DGen routes share bounds and footprint identities,
whereas the CityEngine export lacks the common OSM identifier and ChordAtlas
changes the asset representation.  Topology and colour-variation statistics
are therefore compared only between the two Sat3DGen paths.  Across systems,
the evidence supports conclusions about compactness, material organisation,
available control mechanisms, and workflow structure.

\section{Architectural implications of the manifest-driven design}

The main value of the manifests is that they turn hidden assumptions into
explicit system state.  A request can record the intended geographic extent,
required inputs, permitted tiles, completed stages, and reason for failure.
This enables incomplete acquisition or processing to stop visibly instead of
leaving an output directory that appears complete.  Exact tile selection is
especially important in a shared cache, where recursively consuming every
available mesh can combine unrelated locations while still producing
syntactically valid output.

However, a manifest is authoritative only when it binds outputs to their full
provenance.  Imagery, coordinate reference system, model checkpoint,
inference options, processing configuration, code revision, and output hashes
all influence the meaning of a cached result.  The current design captures
several of these fields but does not yet enforce all of them consistently.
Stable geometric identity is therefore useful for request reuse, but it must be
paired with configuration and source identity before cache reuse can imply
reproducibility.

The cross-language boundary introduces a related transaction problem.  Python
can finish publication before Java has converted and loaded the corresponding
asset. A stronger design would represent processing, conversion, and
visible loading as separate acknowledged stages and would commit the final
state only after the user confirms success.

\section{Implications of the reliability findings}

The large-image comparison shows that nominal overlap alone does not guarantee
a continuous surface.  With hard cropping and equal weights, the set of
contributing windows changes abruptly at regular intervals, turning local
prediction disagreement into visible height steps.  Fractional placement and
raised-cosine weighting substantially reduce those steps without increasing
the measured GPU peak, at the cost of modest additional runtime.  Feathering
smooths disagreement rather than resolving its cause; broad height bias shared
by a window can therefore remain as a gradual surface variation.

The data audit shows that structural completeness and semantic validity are
different properties.  Files may be present while containing duplicates,
cross-partition leakage, inconsistent coordinates, or candidates outside their
intended spatial support.  Dataset preparation should therefore validate
identity, containment, and geographic separation in addition to file counts
and schema shape.  The London workflow is useful as prepared project data, but
these findings prevent treating it as a validated training benchmark.

The mesh experiments reveal the same distinction between execution and
correctness.  A deterministic merge can still select non-contiguous inputs;
a configured vertical threshold can remain ineffective if it is not applied
by the matching predicate; and a DSM correction can produce a complete file
while displacing surfaces that should not be treated as terrain.  The relevant
engineering principle is that every geometric assumption should have both an
executable invariant and a regression case that can falsify it.

Building extraction also creates a trade-off between modularity and geometric
quality.  Footprint association produces independently addressable assets,
which is valuable for editing and procedural processing, but cropping and
capping do not by themselves guarantee closed, manifold, or semantically
separated buildings.  Topology and surface accuracy therefore require explicit
post-processing and evaluation rather than being inferred from successful OBJ
publication.

These observations lead to four general requirements for reliable integration:

\begin{itemize}
  \item validate spatial and semantic relationships, not only the presence and
  shape of files;
  \item make tile membership and provenance authoritative at every stage;
  \item test every geometric assumption and specify how conflicting attributes
  are combined.
\end{itemize}

\section{Contribution assessment}

The contribution is a hybrid reconstruction and integration method rather than
a replacement for either learned reconstruction or procedural modelling.  It
extends the released regional wrapper with continuous density and colour
fusion, then connects the resulting surface to a shared spatial frame,
building identity, external geographic context, modular publication, and
executable data-transfer contracts.  The London data analysis and controlled
geometry tests expose assumptions that would otherwise remain hidden at these
interfaces.

The project builds on Sat3DGen, ChordAtlas, ProceX, BigSUR, and FrankenGAN
\cite{qian2026sat3dgen,kellyChordAtlas2024,kelly2011procex,
kelly2017bigsur,kelly2018frankengan}.  The project-specific contribution is the
London data workflow, large-image density and colour method, geospatial
mesh-processing method, ChordAtlas bridge, and evaluation developed around
these foundations.

\section{Threats to validity}

\subsection{Construct validity}

Most reported measures concern structure and process: coordinates, manifests,
vertices, faces, boundary discontinuity, topology, cache state, and publication
state.  These measures show whether transformations and contracts operate as
specified, but they are not substitutes for geometric or perceptual accuracy.
In particular, smoother seams, complete colour coverage, footprint overlap,
and successful publication do not by themselves establish fa\c{c}ade fidelity
or visual realism.

\subsection{Internal validity and provenance}

The large-image records bind the input plan, script hashes, processing
parameters, and output hashes, but identify the model by repository name rather
than one immutable revision shared explicitly by the density and colour
stages.  The current deployment also resolves those scripts through configured
local paths.  Packaging the scripts with the bridge and binding both stages to
one model revision would make the experiment portable.  At the consumer
boundary, Python requires RGB for the large-image contract, whereas Java can
still fall back from colour parsing; the same requirement should be enforced
at both sides of the transaction.

\subsection{External and statistical validity}

The evaluation is concentrated in one London workflow.  The fusion comparison
uses a matched input, but city, morphology, acquisition, and checkpoint were
not varied.  External services also introduce network, credential, and
licensing dependencies.  Broader claims require a versioned multi-area
evaluation with independent geometry and appearance references.

\section{Priorities for further work}

Three experiments follow directly from the findings.  First, the matched
fusion study should be repeated across multiple urban morphologies and model
revisions, with independent surface and street-level colour references.  A
release-default hard-box run, the fractional-feather method, and relevant
weighting ablations would separate boundary continuity from metric and
perceptual accuracy.

Secondly, the global density path and the independent-mesh path should be run
over the same area and coordinate contract.  Revised vertical stitching and
DSM correction can then be assessed stage by stage for topology, height error,
and surface preservation rather than compared through unrelated cached inputs.

Thirdly, the model revision, scripts, RGB requirement, MiniMesh conversion,
and visible loading acknowledgement should become one portable publication
contract.  The resulting assets can then support a matched CityEngine/CGA
authoring study that extends the present representation example with shared
identity, inputs, and edit tasks, followed by a multi-user Ubiq study of
interactive selection and editing.

%%%%%%%%%%%%%%%%%%%%%%%%%%%%%%%%%%%%%%%%%%%%%%%%%%%%%%%%%%%%%%%%%%%%%%%%%%%%%%%%
\chapter{Conclusion} \label{Chap6}\label{chap:conclusion}

\section{Research outcome}

This dissertation asked to what extent overlapping satellite-conditioned
predictions can be combined into a coloured regional mesh and integrated with
geospatial context in ChordAtlas.  The study establishes both parts of this
workflow.  A rectangular large-image method fuses density predictions at
sub-voxel positions, replaces abrupt overlap changes with smooth weights, and
recovers colour on the fixed global surface.  The resulting mesh is then
placed in a shared spatial frame, associated with OSM footprints and DSM
context, separated into modular assets, and published to ChordAtlas.

The matched London experiment showed substantially lower window-boundary
discontinuity without increasing the measured GPU-memory peak, and the colour
pass covered the complete surface without changing its geometry.  The current
on-demand integration retained that colour through coordinate conversion, DSM
processing, footprint extraction, Java parsing, and building publication.
Manifests, stable selection identities, cache checks, and explicit stage states
make the process inspectable across the Python--Java boundary.

\section{Principal contributions}

The work makes three main contributions:

\begin{enumerate}
  \item It extends large-image Sat3DGen inference to rectangular mosaics through
  sub-voxel density placement and smooth overlap fusion, and recovers global
  vertex colour in a memory-bounded second pass without altering the extracted
  geometry.
  \item It defines a coordinate-explicit path from fused or independently
  assembled geometry to footprint-associated ChordAtlas assets, preserving
  selection identity and vertex appearance while adding OSM and DSM context.
  \item It introduces and evaluates manifest, cache, validation, and
  publication contracts across Python and Java, using data audits, matched
  fusion measurements, coordinate checks, topology inspection, and controlled
  regression cases.
\end{enumerate}

These contributions provide a practical connection between learned urban
reconstruction and procedural authoring.  Sat3DGen supplies site-conditioned
density and appearance, while OSM, DSM, and ChordAtlas add spatial identity,
height context, modular asset boundaries, and an environment for further
editing.  CityEngine clarifies the complementary value of rule-driven
semantics and controllability, while the retained ChordAtlas exports illustrate
the representation difference between per-vertex colour and structured
procedural material maps.

\section{Implications and next steps}

The evaluation identifies the next research priorities.  Multi-area tests with
independent geometry and appearance references are needed to determine whether
smoother window transitions also improve reconstruction accuracy.  The
independent-mesh path requires corrected vertical stitching and more selective
DSM processing, while the published building assets require topology repair.
For deployment, density and colour should be bound to one immutable model
revision, the large-image scripts should be packaged with the bridge, and
readiness should include Java conversion and visible loading acknowledgement.

The current CityEngine export and retained ChordAtlas representation examples
can be extended into a controlled authoring study with shared identity, inputs,
and edit tasks, while
Ubiq can test interactive and collaborative use of the resulting assets.

\section{Closing statement}

The dissertation connects bounded satellite predictions to a continuous
coloured surface and an auditable urban-asset workflow.  Its central value lies
in combining learned
density and appearance with explicit coordinates, geographic identity, and
publication state, so that generated geometry can participate in a larger
editable system.  This provides a concrete foundation for evaluating regional
reconstruction, procedural authoring, and collaborative XR with the same
spatially identified assets.

%%%%%%%%%%%%%%%%%%%%%%%%%%%%%%%%%%%%%%%%%%%%%%%%%%%%%%%%%%%%%%%%%%%%%%%%%%%%%%%%
\renewcommand{\bibname}{References}
\bibliographystyle{ieeetr}
%\bibliography{Bibliography.bib}
\bibliography{references.bib}
%%%%%%%%%%%%%%%%%%%%%%%%%%%%%%%%%%%%%%%%%%%%%%%%%%%%%%%%%%%%%%%%%%%%%%%%%%%%%%%%
% APPENDIX
\begin{appendices}
\chapter{Implementation and Reproducibility Record}\label{app:evidence-record}

\section{Evaluated components}

The empirical study uses four principal implementations: the London
data-preparation workflow, the Sat3DGen large-image extension, the active mesh
assembly package, and the ChordAtlas integration layer.
The project source code and submission package are available at
\url{https://github.com/renjie22122-maker/FinalSubmit}.
Runtime import and entry-point tracing identify the implementation used for
each experiment.  The published Sat3DGen,
ChordAtlas, BigSUR, FrankenGAN, CityEngine, and Ubiq systems are referenced in
the chapters where their methods or interfaces are discussed.

\section{Component-test results}

Table~\ref{tab:appendix-component-tests} records the current component-level
verification results.  The Python suite exercises data prerequisites,
coordinates, tile planning, geometry processing, manifests, cache behaviour,
per-building publication, appearance references, and rollback.  The Java
validators exercise coordinate conversion, selection and cache contracts, OBJ
slicing, vertex-colour transfer, MiniMesh publication, and appearance-bridge
components.

\begin{table}[H]
\centering
\caption{Detailed component-verification results.}
\label{tab:appendix-component-tests}
\small
\begin{tabularx}{\textwidth}{p{0.25\textwidth}p{0.18\textwidth}X}
\toprule
Test group & Result & Principal properties exercised\\
\midrule
Python integration suite & 124/124 passed & Input validation, coordinates,
geometry rules, manifests, cache state, publication, rollback, roof references,
panorama acquisition, and large-image bridge contracts.\\
Large-image algorithm regressions & 9/9 passed & Fractional placement,
coverage, feather continuity, colour-coordinate mapping, weighted RGB, and PLY
round trip.\\
Java bridge validators & 8/8 passed & MiniMesh conversion, selection identity,
OBJ slicing, vertex colour, coordinate transfer, appearance integration, and
GUI selection dispatch.\\
Controlled geometry cases & 2 counterexamples reproduced & Bottom-surface
classification and the vertical constraint in vertex stitching.\\
\bottomrule
\end{tabularx}
\end{table}

The controlled geometry cases correspond to the predicate results discussed in
Chapters~\ref{Chap4} and~\ref{Chap5}.  Quantitative figures were generated from
the same evaluated artefacts and record their relevant populations in their
captions or surrounding text.

\section{Selected implementation excerpts}

The following short excerpts connect Algorithms~\ref{alg:fractional-density-fusion}
and~\ref{alg:vertex-colour-recovery} to the submitted Python implementation.
They are included verbatim from the project-authored scripts; the complete
files, command-line interfaces, validation logic, and tests remain under
\path{research/scripts}.

\lstinputlisting[
  language=Python,
  firstline=157,
  lastline=188,
  caption={Raised-cosine axis weights used by the large-image density-fusion script.},
  label={lst:raised-cosine-weights}
]{research/scripts/run_sat3dgen_big_image_app192.py}

\lstinputlisting[
  language=Python,
  firstline=241,
  lastline=274,
  caption={Exact support-weight selection used by the memory-bounded colour pass.},
  label={lst:window-vertex-weights}
]{research/scripts/colorize_sat3dgen_big_mesh.py}

\section{Model-example comparison record}

The representation comparison is generated from the retained
\texttt{modelsExample} directory.  Only the \texttt{BigImageModel} and
\texttt{satelliteImageModel} request/result packages provide directly matched
identity: both name OSM way~369245408 and contain the same target bounds and
footprint geometry.  Their building OBJ files provide the geometry, topology,
and vertex-colour evidence, while the result manifests provide the recorded DSM
status and coverage reported in
Section~\ref{sec:matched-building-comparison}.

The CityEngine directory retains a CityEngine 2026 export log, OBJ, MTL, and
three texture images, but no common OSM identifier.  The two route-labelled
ChordAtlas directories retain OBJ/MTL outputs and diffuse/normal/specular map
sets; the OBJ files do not embed the footprint identity.  The standalone
ChordAtlas directory has no matching location or run manifest and is used only
as an illustrative reference.  The reproducible audit and software renderer
are stored in \texttt{research/scripts}; their JSON/CSV outputs are stored in
\texttt{research/results/models\_example}.  The rendered comparison figures use
derived vertex colour or stored material maps and do not embed the source
satellite mosaic.

\section{Data and licence notes}

OpenStreetMap-derived content requires Open Database Licence attribution
\cite{openstreetmapData}.  The Environment Agency 2022 first-return DSM is
used under the Open Government Licence and retains its stated coordinate and
height reference \cite{environmentAgencyDSM2022,nationalArchivesOGL3}.  Google
imagery is not redistributed in the generated audit figures.  Credentials
found in local prototypes are omitted from the report.


\chapter{Additional Matched London Building Comparison}\label{app:models-example-2}

\section{Matched satellite-derived outputs}

The \texttt{modelsExample2} record provides a second geographically matched
case for the large-image and independent-tile routes.  The two request
packages contain identical selection bounds and footprint geometry for OSM
ways~33719501 and~1455280501.  The large-image route processes a contiguous
4-by-4 satellite mosaic through 100 overlapping inference windows, whereas
the independent route assembles 20 separately generated meshes.  Both result
packages record successful DSM sampling for all vertices in the extracted
building assets.

Table~\ref{tab:models-example-2-metrics} records the building-level geometry
reported for both requested footprints.  Vertical relief is the range within
each saved model and is not a surveyed building-height measurement.

\begin{table}[H]
\centering
\caption{Recorded geometry metrics for the matched \texttt{modelsExample2}
building assets.}
\label{tab:models-example-2-metrics}
\small
\begin{tabularx}{\textwidth}{>{\raggedright\arraybackslash}p{0.18\textwidth}>{\raggedright\arraybackslash}p{0.14\textwidth}rrr>{\raggedright\arraybackslash}X}
\toprule
Route & OSM way & Vertices & Faces & Relief (m) & Boundary / non-manifold
edges\\
\midrule
Large-image & 33719501 & 122,119 & 242,601 & 37.75 & 1,617 / 0\\
Independent tiles & 33719501 & 126,983 & 259,125 & 40.22 & 4,807 / 6,072\\
Large-image & 1455280501 & 561 & 1,022 & 16.73 & 98 / 0\\
Independent tiles & 1455280501 & 300 & 578 & 19.68 & 111 / 58\\
\bottomrule
\end{tabularx}
\end{table}

For the primary footprint, the large-image asset has fewer boundary and
non-manifold edges than the independent-tile asset in this case.  All four
extracted surfaces remain open rather than watertight.  The comparison
therefore extends the case evidence in Chapter~\ref{Chap4} without treating
mesh counts as a measure of reconstruction accuracy.

Figure~\ref{fig:models-example-2-matched} compares the primary footprint under
shared cameras and projection bounds.  The renders use the RGB attributes
stored by each satellite-derived route; no source satellite image is embedded
in the figure.  The generating script and machine-readable audit are retained
under \texttt{research/scripts} and
\texttt{research/results/models\_example2}, respectively.

\begin{figure}[!htbp]
\centering
\includegraphics[width=\textwidth]{models_example2_matched_sat3dgen.pdf}
\caption[Additional matched Sat3DGen example.]{Additional matched comparison
for OSM way~33719501.  Panels (a,c) show the large-image output and panels
(b,d) the independent-tile output under shared oblique and plan-view bounds.
Colours are the per-vertex RGB attributes stored by the respective
satellite-derived routes.}
\label{fig:models-example-2-matched}
\end{figure}
\FloatBarrier

\section{Six representations of the selected building}

The \texttt{modelsExample2} record contains six representations of the same
selected London Trocadero building.  They comprise the large-image and
independent-tile Sat3DGen models, the corresponding ChordAtlas outputs, the
CityEngine model, and the original ChordAtlas model.  The two Sat3DGen result
manifests identify the building as OSM way~33719501.  The later OBJ/MTL exports
belong to the same selected case, although they do not repeat the OSM
identifier inside the model files.

The two Sat3DGen routes retain dense, irregular satellite-conditioned surfaces
and encode appearance as per-vertex colour.  Their corresponding ChordAtlas
outputs replace vertex colour with UV-mapped diffuse, normal, and specular
material structures while maintaining a similar horizontal building extent.
The CityEngine model provides the most compact and regular envelope; its doors
and windows are represented in shared image materials rather than as separate
openings.  The ChordAtlas packages contain independently materialised facade
regions, including window-labelled materials and, in the independent-tile
output, door-labelled materials.  These differences describe geometry and
appearance representation; greater face or material counts do not alone imply
greater reconstruction accuracy.

Table~\ref{tab:models-example-2-six-packages} describes the stored geometry
and appearance representation in every package.  Face counts are reported as
triangles after OBJ triangulation, except that the CityEngine row also retains
its original polygon-record count.

\begin{table}[H]
\centering
\caption[Six-package geometry and appearance records.]{Geometry and appearance
records for the six representations in \texttt{modelsExample2}.}
\label{tab:models-example-2-six-packages}
\scriptsize
\begin{tabularx}{\textwidth}{>{\raggedright\arraybackslash}p{0.18\textwidth}>{\raggedright\arraybackslash}p{0.20\textwidth}>{\raggedright\arraybackslash}p{0.34\textwidth}>{\raggedright\arraybackslash}X}
\toprule
Representation & OBJ geometry & Appearance encoding & Method or stage\\
\midrule
Large-image Sat3DGen & 122,119 vertices; 242,601 triangles & RGB on every
vertex; no UV material & Large-image route output\\
Large-image ChordAtlas & 89,647 vertices; 728,314 triangles & 69,041 UV
records; 527 materials; 394 diffuse--normal--specular triplets &
ChordAtlas output from the large-image route\\
CityEngine 2026 & 66 vertices; 23 polygon records (64 triangles after
triangulation) & 52 UV records; two materials and two diffuse images &
CityEngine model\\
Independent-tile Sat3DGen & 126,983 vertices; 259,125 triangles & RGB on every
vertex; no UV material & Independent-tile route output\\
Independent-tile ChordAtlas & 226,056 vertices; 2,597,150 triangles & 156,686
UV records; 3,151 materials; 2,841 diffuse--normal--specular triplets &
ChordAtlas output from the independent-tile route\\
Original ChordAtlas & 149,021 vertices; 1,027,970 triangles & 125,922 UV
records; 572 materials; 424 diffuse--normal--specular triplets &
Original ChordAtlas model\\
\bottomrule
\end{tabularx}
\end{table}

Figure~\ref{fig:models-example-2-six-package} presents the six outputs with a
common oblique camera direction and metric display scale.  Each model is
translated to the centre of its panel without rescaling; the original
ChordAtlas model is rotated by 180 degrees about its vertical axis so that the
same side of the building is shown.

\begin{figure}[!htbp]
\centering
\includegraphics[width=\textwidth]{models_example2_six_method_comparison.pdf}
\caption[Six representations of the same selected London building.]{Six
representations of the same selected London Trocadero building.  Panels (a,b)
show the large-image Sat3DGen output before and after ChordAtlas processing;
panels (d,e) show the corresponding independent-tile pair; panel (c) shows
the CityEngine model; and panel (f) shows the original ChordAtlas model.  The
two Sat3DGen meshes display stored vertex RGB, whereas the CityEngine and
ChordAtlas models display UV-mapped diffuse materials.  All panels use a
common camera direction and metric scale; models are centred without
rescaling, and panel (f) uses the stated presentation rotation.}
\label{fig:models-example-2-six-package}
\end{figure}

\FloatBarrier

\end{appendices}
%%%%%%%%%%%%%%%%%%%%%%%%%%%%%%%%%%%%%%%%%%%%%%%%%%%%%%%%%%%%%%%%%%%%%%%%%%%%%%%%
% GLOSSARY
\clearpage
\printglossaries

% INDEX?

\end{document}
